
\chapter{Étale cohomology: the starting points}\label{I}
\blfootnote{by P.\ Deligne, written up by J.F.\ Boutot}

\section{Grothendieck topologies}\label{I:1}

Grothendieck topologies first appeared as the foundation of Grothendieck's
theory of descent (cf. SGA 1 VI, VIII); their use in the corresponding
cohomology theories came later. We follow the same path here: by formalizing
the classical notions of localization, local property, and patching
(\ref{I:1-1}, \ref{I:1-2}, \ref{I:1-3}), we obtain the general concept of a
Grothendieck topology (\ref{I:1-6}); to justify its introduction into
algebraic geometry, we prove a faithfully flat descent theorem
(\ref{I:1-4}), generalizing the classical theorem 90 of Hilbert
(\ref{I:1-5}).

The reader will find a more complete, but concise, exposition of the
formalism in Giraud \cite{gi64}. M. Artin's notes, “Grothendieck
topologies” \cite{ar62} (Chapters I to III), also remain useful. The
866 pages of Exposés I to V of SGA 4 are valuable when one considers exotic
topologies, such as the one giving rise to crystalline cohomology; in order
to use the étale topology, which is so close to classical intuition, it is
not essential to read them.

\subsection{Sieves}\label{I:1-1}

Let $X$ be a topological space and let $f:X\to \dR$ be a real-valued
function on $X$. The continuity of $f$ is a local property; in other words,
if $f$ is continuous on every sufficiently small open subset of $X$, then
$f$ is continuous on all of $X$. To formalize the notion of a “local
property,” we introduce some definitions.

A set $\cU$ of open subsets of $X$ is called a \emph{sieve} if, for every
$U\in\cU$ and $V\subset U$, one has $V\in\cU$. A sieve is called
\emph{covering} if the union of all the open subsets belonging to it is
equal to $X$.

Given a family $\{U_i\}$ of open subsets of $X$, the sieve generated by
$\{U_i\}$ is, by definition, the set of open subsets $U$ of $X$ such that
$U$ is contained in one of the $U_i$.

A property $P(U)$, defined for every open subset $U$ of $X$, is called
\emph{local} if, for every covering sieve $\cU$ of every open subset $U$
of $X$, $P(U)$ is true if and only if $P(V)$ is true for every
$V\in \cU$. For example, given $f:X\to \dR$, the property “$f$ is
continuous on $U$” is local.

\subsection{Sheaves}\label{I:1-2}

Let us make precise the notion of a function given locally on $X$.

\subsubsection{The sieve point of view}\label{I:1-2-1}

Let $\cU$ be a sieve of open subsets of $X$. A function given
$\cU$-locally on $X$ consists of a function $f_U$ on $U$ for every
$U\in \cU$, such that if $V\subset U$ then $f_V=f_U|V$.

\subsubsection{The Čech point of view}\label{I:1-2-2}

If the sieve $\cU$ is generated by a family of open subsets $U_i$ of $X$,
then giving a function $\cU$-locally amounts to giving a function $f_i$
on each $U_i$ such that
$f_i|{U_i\cap U_j} = f_j|{U_i\cap U_j}$.

In other words, if $Z=\coprod U_i$, then giving a function
$\cU$-locally amounts to giving a function on $Z$ that is constant on the
fibres of the natural projection $Z\to X$.

\subsubsection{}\label{I:1-2-3}

The continuous functions form a sheaf. This means that for every covering
sieve $\cU$ of an open subset $V$ of $X$, and every function
$\{f_U\}$ given $\cU$-locally such that each $f_U$ is continuous on
$U$, there exists a unique continuous function $f$ on $V$ such that
$f|U=f_U$ for every $U\in \cU$.

\subsection{Stacks}\label{I:1-3}

Let us now make precise the notion of a vector bundle given locally on $X$.

\subsubsection{The sieve point of view}\label{I:1-3-1}

Let $\cU$ be a sieve of open subsets of $X$. A vector bundle given
$\cU$-locally on $X$ consists of
\begin{enumerate}[\indent a)]
  \item a vector bundle $E_U$ on each $U\in \cU$,
  \item if $V\subset U$, an isomorphism
    $\rho_{U,V}:E_V\iso E_U|V$, satisfying
  \item if $W\subset V\subset U$, the diagram
    \[\xymatrix{
      E_W \ar[r]^-{\rho_{U,W}} \ar[dr]_-{\rho_{V,W}}
        & E_U|W \\
      & E_V|W \ar[u]_-{\rho_{U,V}|W}
    }\]
    commutes, that is,
    $\rho_{U,W} = (\mbox{$\rho_{U,V}$ restricted to $W$}) \circ \rho_{V,W}$.
\end{enumerate}

\subsubsection{The Čech point of view}\label{I:1-3-2}

If the sieve $\cU$ is generated by a family of open subsets $U_i$ of
$X$, then giving a vector bundle $\cU$-locally amounts to giving:
\begin{enumerate}[\indent a)]
  \item a vector bundle $E_i$ on each $U_i$,
  \item if $U_{ij}=U_i\cap U_j=U_i\times_X U_j$, an isomorphism
    $\rho_{ji}:E_i|U_{ij}\iso E_j|U_{ij}$, such that
  \item if $U_{ijk}=U_i\times_X U_j\times_X U_k$, the diagram
    \[\xymatrix{
      E_i |U_{ijk} \ar[r]^-{\rho_{ki}|U_{ijk}} \ar[dr]_-{\rho_{ji}|U_{ijk}}
        & E_k|U_{ijk} \\
      & E_j|U_{ijk} \ar[u]_-{\rho_{kj}|U_{ijk}}
    }\]
    commutes, that is, $\rho_{ki}=\rho_{kj}\circ\rho_{ji}$ on
    $U_{ijk}$.
\end{enumerate}

In other words, if $Z=\coprod U_i$ and $\pi:Z\to X$ is the natural
projection, then giving a vector bundle $\cU$-locally amounts to giving:
\begin{enumerate}[\indent a)]
  \item a vector bundle $E$ on $Z$,
  \item if $x$ and $y$ are two points of $Z$ such that
    $\pi(x)=\pi(y)$, an isomorphism $\rho_{yx}:E_x\iso E_y$ between the
    fibres of $E$ at $x$ and $y$, depending continuously on $(x,y)$,
    and such that
  \item if $x$, $y$, and $z$ are three points of $Z$ such that
    $\pi(x)=\pi(y)=\pi(z)$, then
    $\rho_{zx}=\rho_{zy}\circ\rho_{yx}$.
\end{enumerate}

\subsubsection{}\label{I:1-3-3}

A vector bundle $E$ on $X$ defines a vector bundle $E_\cU$ given
$\cU$-locally, namely the system of restrictions $E_U$ of $E$ to the
objects of $\cU$. The fact that the notion of a vector bundle is local can
be expressed as follows: for every covering sieve $\cU$ of $X$, the
functor $E\mapsto E_\cU$ from vector bundles on $X$ to vector bundles
given $\cU$-locally is an equivalence of categories.

\subsubsection{}\label{I:1-3-4}

If in \ref{I:1} one replaces “open subset of $X$” by “subset of $X$,”
one obtains the notion of a sieve of subspaces of $X$. In this setting
there are also patching theorems. For example, let $X$ be a space and
let $\cC$ be a sieve of subspaces of $X$ generated by a locally finite
closed covering of $X$. Then the functor $E\mapsto E_\cC$ from vector
bundles on $X$ to vector bundles given $\cC$-locally is an equivalence
of categories.

In algebraic geometry, it is also useful to consider “sieves of spaces
over $X$”; this is what we shall see in the next paragraph.

\subsection{Faithfully flat descent}\label{I:1-4}

\subsubsection{}\label{I:1-4-1}

In the setting of schemes, the Zariski topology is not fine enough for the
study of nonlinear problems, and one is led to replace the open immersions
in the preceding definitions by more general morphisms. From this point of
view, descent techniques appear as localization techniques. Thus the
following descent statement can be expressed by saying that the properties
under consideration are local for the faithfully flat topology (a morphism
of schemes is called faithfully flat if it is flat and surjective).

\begin{proposition}\label{I:1-4-2}
Let $A$ be a ring and $B$ a faithfully flat $A$-algebra. Then:
\begin{enumerate}[(i)]
  \item A sequence $\Sigma=(M'\to M\to M'')$ of $A$-modules is exact if
    the sequence $\Sigma_{(B)}$ deduced from it by extension of scalars to
    $B$ is exact.
  \item An $A$-module $M$ is of finite type (resp. of finite presentation,
    flat, locally free of finite rank, invertible (i.e. locally free of rank
    one)) if the $B$-module $M_{(B)}$ has the corresponding property.
\end{enumerate}
\end{proposition}
\begin{proof}
(i) Since the functor $M\mapsto M_{(B)}$ is exact (by flatness of $B$),
it suffices to show that if an $A$-module $N$ is nonzero, then
$N_{(B)}$ is nonzero. If $N$ is nonzero, $N$ contains a nonzero cyclic
submodule $A/\fa$; hence $N_{(B)}$ contains the cyclic submodule
$(A/\fa)_{(B)}=B/\fa B$, which is nonzero by surjectivity of the
structural morphism $\varphi:\spec(B)\to\spec(A)$ (if $V(\fa)$ is
nonempty, then $\varphi^{-1}(V(\fa))=V(\fa B)$ is nonempty).

(ii) For every family $(x_i)$ of elements of $M_{(B)}$, there exists a
finitely generated submodule $M'$ of $M$ such that $M'_{(B)}$ contains
the $x_i$. If $M_{(B)}$ is of finite type and the $x_i$ generate
$M_{(B)}$, then $M'_{(B)}=M_{(B)}$, and hence $M'=M$ and $M$ is of
finite type.
\end{proof}

If $M_{(B)}$ is of finite presentation, the preceding argument yields a
surjection $A^n\to M$. If $N$ is the kernel of this surjection, then the
$B$-module $N_{(B)}$ is of finite type, so $N$ is of finite type and
$M$ is of finite presentation. The assertion for “flat” follows at once
from (i); “locally free of finite rank” means “flat and of finite
presentation,” and the rank can be tested after extension of scalars to
fields.

\subsubsection{}\label{I:1-4-3}

Let $X$ be a scheme and let $\cS$ be a class of $X$-schemes stable under
fibre products over $X$. A class $\cU\subset\cS$ is a \emph{sieve} on
$X$ (relative to $\cS$) if, for every morphism $\varphi:V\to U$ of
$X$-schemes with $U,V\in\cS$ and $U\in\cU$, one has $V\in\cU$.
The sieve \emph{generated} by a family $\{U_i\}$ of $X$-schemes in
$\cS$ is the class of $V\in\cS$ for which there exists a morphism of
$X$-schemes from $V$ to one of the $U_i$.

\subsubsection{}\label{I:1-4-4}

Let $\cU$ be a sieve on $X$. A quasi-coherent module given
$\cU$-locally on $X$ consists of
\begin{enumerate}[\indent a)]
  \item a quasi-coherent module $E_U$ on every $U\in\cU$,
  \item for every $U\in\cU$ and every morphism $\varphi:V\to U$ of
    $X$-schemes in $\cS$, an isomorphism
    $\rho_\varphi:E_V\iso\varphi^*E_U$, such that
  \item if $\psi:W\to V$ is a morphism of $X$-schemes in $\cS$, the
    diagram
    \[\xymatrix{
      E_W \ar[r]^-{\rho_{\varphi\circ\psi}} \ar[dr]_-{\rho_\psi}
        & \psi^*\varphi^* E_U \\
      & \psi^* E_V \ar[u]_-{\psi^*\rho_\varphi}
    }\]
    commutes, that is,
    $\rho_{\varphi\circ\psi}=(\psi^*\rho_\varphi)\circ\rho_\psi$.
\end{enumerate}

If $E$ is a quasi-coherent module on $X$, we write $E_\cU$ for the
module given $\cU$-locally whose value is $\varphi_U^*E$ on
$\varphi:U\to X$ and such that, for every morphism $\psi:V\to U$, the
restriction isomorphism $\rho_\psi$ is the canonical isomorphism
$E_V=(\varphi_U\circ\psi)^*E\iso\psi^*\varphi_U^*E=\psi^*E_U$.

\begin{theorem}\label{I:1-4-5}
Let $\{U_i\}\in\cS$ be a finite family of $X$-schemes flat over $X$
such that $X$ is the union of the images of the $U_i$, and let $\cU$
be the sieve generated by $\{U_i\}$. Then the functor
$E\mapsto E_\cU$ is an equivalence from the category of quasi-coherent
modules on $X$ to the category of quasi-coherent modules given
$\cU$-locally.
\end{theorem}
\begin{proof}
We treat only the case in which $X$ is affine and $\cU$ is generated by
an affine $X$-scheme $U$ faithfully flat over $X$. The reduction to this
case is formal. Put $X=\spec(A)$ and $U=\spec(B)$.

If $U\to X$ admits a section, then $X$ belongs to the sieve $\cU$ and
the assertion is clear. We reduce to this case.

A quasi-coherent module given $\cU$-locally defines modules $M'$,
$M''$, and $M'''$ on $U$, $U\times_XU$, and
$U\times_XU\times_XU$, and isomorphisms
$\rho:p^*M^\bullet\simeq M^\bullet$ for every projection morphism $p$
between these spaces; this is a \emph{cartesian diagram}
\[\xymatrix{
  M^* : M' \ar@<2pt>[r] \ar@<-2pt>[r]
    & M'' \ar[r] \ar@<4pt>[r] \ar@<-4pt>[r]
    & M'''
}\]
over
\[\xymatrix{
  U_* : U
    & U\times_X U \ar@<2pt>[l] \ar@<-2pt>[l]
    & U\times_X U\times_X U \ar[l] \ar@<4pt>[l] \ar@<-4pt>[l] \mbox{.}
}\]
Conversely, $M^*$ determines the module given $\cU$-locally: for
$V\in\cU$ there exists $\varphi:V\to U$, and we put
$M_V=\varphi^*M'$; for $\varphi_1,\varphi_2:V\to U$ there is a natural
identification
$\varphi_1^*M'\simeq(\varphi_1\times\varphi_2)^*M''\simeq\varphi_2^*M'$,
and by using $M'''$ one sees that these identifications are compatible,
so the definition is legitimate. In short, giving a module
$\cU$-locally is equivalent to giving a cartesian diagram $M^*$ over
$U_*$.

In algebraic terms, giving $M^*$ amounts to giving a cartesian diagram of
modules
\[\xymatrix{
  M' \ar@<3pt>[r]|-{\partial_0} \ar@<-3pt>[r]|-{\partial_1}
    & M'' \ar@<6pt>[r]|-{\partial_0} \ar[r]|-{\partial_1} \ar@<-6pt>[r]|-{\partial_2}
    & M'''
}\]
over the diagram of rings
\[\xymatrix{
  B \ar@<3pt>[r]|-{\partial_0} \ar@<-3pt>[r]|-{\partial_1}
    & B\otimes_A B \ar@<6pt>[r]|-{\partial_0} \ar[r]|-{\partial_1} \ar@<-6pt>[r]|-{\partial_2}
    & B\otimes_A B\otimes_A B
}\]
(more precisely, $\partial_i(bm)=\partial_i(b)\cdot\partial_i(m)$, the
usual identities such as $\partial_0\partial_1=\partial_0\partial_0$
hold, and “cartesian” means that the morphisms
$\partial_i:M'\otimes_{B,\partial_i}(B\otimes_A B)\to M''$ and
$M''\otimes_{B\otimes_A B,\partial_i}(B\otimes_A B\otimes_A B)\to M'''$
are isomorphisms).

The functor $E\mapsto E_\cU$ becomes the functor that associates to an
$A$-module $M$ the diagram
\[\xymatrix{
  M^* = (M\otimes_A B \ar@<2pt>[r] \ar@<-2pt>[r]
    & M\otimes_A B\otimes_A B \ar@<-4pt>[r] \ar[r] \ar@<4pt>[r]
    & M\otimes_A B\otimes_A B\otimes_A B)\text{.}
}\]
Its right adjoint is the functor
\[\xymatrix{
  (M' \ar@<2pt>[r] \ar@<-2pt>[r]
    & M'' \ar@<-4pt>[r] \ar[r] \ar@<4pt>[r]
    & M''') \ar@{|->}[r]
    & \ker(M' \ar@<2pt>[r] \ar@<-2pt>[r]
    & M''')\text{.}
}\]
We must prove that the adjunction maps
\[
  M \to\ker(M\otimes_A B\rightrightarrows M\otimes_A B\otimes_A B)
\]
and
\[
  \ker(M'\rightrightarrows M'')\otimes_A B \to M'
\]
are isomorphisms. By (\ref{I:1-4-2}.i), it suffices to prove this after
a faithfully flat base change $A\to A'$ (with $B$ becoming
$B'=B\otimes_AA'$). Taking $A'=B$ reduces us to the case in which
$U\to X$ admits a section.
\end{proof}

\subsection{A special case: Hilbert's theorem 90}\label{I:1-5}

\subsubsection{}\label{I:1-5-1}

Let $k$ be a field, let $k'$ be a Galois extension of $k$, and let
$G=\gal(k'/k)$. Then the homomorphism
\begin{align*}
  k'\otimes_k k' &\to \oplus_{\sigma\in G} k' \\
  x\otimes y     &\mapsto \{x\cdot \sigma(y)\}_{\sigma\in G}
\end{align*}
is bijective.

It follows that giving a module locally for the sieve on $\spec(k)$
generated by $\spec(k')$ is equivalent to giving a $k'$-vector space
equipped with a semilinear action of $G$, that is:
\begin{enumerate}[\indent a)]
  \item a $k'$-vector space $V'$,
  \item for every $\sigma\in G$, an endomorphism $\varphi_\sigma$ of the
    underlying group of $V'$ such that
    $\varphi_\sigma(\lambda v)=\sigma(\lambda)\varphi_\sigma(v)$ for all
    $\lambda\in k'$ and $v\in V'$, satisfying the condition
  \item for all $\sigma,\tau\in G$,
    $\varphi_{\tau\sigma}=\varphi_\tau\circ\varphi_\sigma$.
\end{enumerate}

Let $V={V'}^G$ be the group of invariants under this action of $G$. It
is a $k$-vector space, and by Theorem \ref{I:1-4-5} one has:

\begin{proposition}\label{I:1-5-2}
The inclusion of $V$ in $V'$ defines an isomorphism
$V\otimes_k k'\iso V'$.
\end{proposition}

In particular, if $V'$ has dimension $1$ and $v'\in V'$ is nonzero, then
$\varphi_\sigma$ is determined by the constant
$c(\sigma)\in {k'}^\times$ such that
$\varphi_\sigma(v')=c(\sigma)v'$, and condition c) becomes
\[
  c(\tau\sigma) = c(\tau)\cdot \tau(c(\sigma))\text{.}
\]
By the proposition, there exists a nonzero invariant vector
$v=\mu v'$, $\mu\in {k'}^\times$. Hence, for every $\sigma\in G$,
\[
  c(\sigma) = \mu\cdot \sigma(\mu^{-1})\text{.}
\]
In other words, every $1$-cocycle of $G$ with values in
${k'}^\times$ is a coboundary:

\begin{corollary}\label{I:1-5-3}
One has $\h^1(G,{k'}^\times)=0$.
\end{corollary}

\subsection{Grothendieck topologies}\label{I:1-6}

We now transpose the definitions of the preceding paragraphs to an
abstract setting that encompasses both topological spaces and schemes.

\subsubsection{}\label{I:1-6-1}

Let $\cS$ be a category and let $U$ be an object of $\cS$. A
\emph{sieve} on $U$ is a subset $\cU$ of $\ob(\cS/U)$ such that if
$\varphi:V\to U$ belongs to $\cU$ and $\psi:W\to V$ is a morphism in
$\cS$, then $\varphi\circ\psi:W\to U$ belongs to $\cU$.

If $\{\varphi_i:U_i\to U\}$ is a family of morphisms, the sieve generated
by the $U_i$ is, by definition, the set of morphisms $\varphi:V\to U$
that factor through one of the $\varphi_i$.

If $\cU$ is a sieve on $U$ and $\varphi:V\to U$ is a morphism, the
restriction $\cU_V$ of $\cU$ to $V$ is, by definition, the sieve on
$V$ consisting of the morphisms $\psi:W\to V$ such that
$\varphi\circ\psi:W\to U$ belongs to $\cU$.

\subsubsection{}\label{I:1-6-2}

A \emph{Grothendieck topology} on $\cS$ consists of the choice, for every
object $U$ of $\cS$, of a set $C(U)$ of sieves on $U$, called
covering sieves, such that the following axioms hold:
\begin{enumerate}[\indent a)]
  \item The sieve generated by the identity of $U$ is covering.
  \item If $\cU$ is a covering sieve on $U$ and $V\to U$ is a
    morphism, then the sieve $\cU_V$ is covering.
  \item A locally covering sieve is covering. In other words, if $\cU$
    is a covering sieve on $U$ and $\cU'$ is a sieve on $U$ such that,
    for every $V\to U$ belonging to $\cU$, the sieve $\cU_V'$ is
    covering, then $\cU'$ is covering.
\end{enumerate}

A category equipped with a Grothendieck topology is called a \emph{site}.

\subsubsection{}\label{I:1-6-3}

Given a site $\cS$, a presheaf on $\cS$ is a contravariant functor
$\sF$ from $\cS$ to the category of sets. For every object $U$ of
$\cS$, the elements of $\sF(U)$ are called sections of $\sF$ over
$U$. For every morphism $V\to U$ and every $s\in\sF(U)$, we denote by
$s|V$ ($s$ restricted to $V$) the image of $s$ in $\sF(V)$.

If $\cU$ is a sieve on $U$, a section given $\cU$-locally consists,
for every $V\to U$ belonging to $\cU$, of a section
$s_V\in\sF(V)$ such that, for every morphism $W\to V$,
$s_V|W=s_W$. The presheaf $\sF$ is called a \emph{sheaf} if, for every
object $U$ of $\cS$, every covering sieve $\cU$ on $U$, and every
section $\{s_V\}$ given $\cU$-locally, there exists a unique section
$s\in\sF(U)$ such that $s|V=s_V$ for every $V\to U$ belonging to
$\cU$.

\emph{Abelian sheaves} are defined similarly by replacing the category of
sets with the category of abelian groups. One shows that the category of
abelian sheaves on $\cS$ is an abelian category with enough injectives. A
sequence $\sF\xrightarrow f\sG\xrightarrow g\sH$ of sheaves is exact if,
for every object $U$ of $\cS$ and every $s\in\sG(U)$ such that
$g(s)=0$, there exists locally a $t$ such that $f(t)=s$; i.e. there
exists a covering sieve $\cU$ on $U$ and, for every $V\in\cU$, a
section $t_V$ of $\sF$ on $V$ such that $f(t_V)=s|V$.

\subsubsection{Examples}\label{I:1-6-4}

We have seen two examples above.

\begin{enumerate}[\indent a)]
  \item Let $X$ be a topological space and let $\cS$ be the category
    whose objects are the open subsets of $X$ and whose morphisms are the
    natural inclusions. The Grothendieck topology on $\cS$ corresponding
    to the usual topology of $X$ is the one for which a sieve $\cU$ on
    an open subset $U$ of $X$ is covering if the union of the open
    subsets belonging to the sieve equals $U$. It is clear that the
    category of sheaves on $\cS$ is equivalent to the category of sheaves
    on $X$ in the usual sense.
  \item Let $X$ be a scheme and let $\cS$ be the category of schemes over
    $X$. The fpqc (faithfully flat quasi-compact) topology on $\cS$ is
    the Grothendieck topology for which a sieve on an $X$-scheme $U$ is
    covering if it is generated by a finite family of flat morphisms whose
    images cover $U$.
\end{enumerate}

\subsubsection{Cohomology}\label{I:1-6-5}

We shall always assume that the category $\cS$ has a final object $X$.
The global sections of an abelian sheaf $\sF$, denoted by $\Gamma\sF$ or
$\h^0(X,\sF)$, form the group $\sF(X)$. The functor
$\sF\mapsto\Gamma\sF$ is a left exact functor from the category of
abelian sheaves on $\cS$ to the category of abelian groups; its derived
functors (or satellites) are denoted by $\h^i(X,-)$. These cohomology
groups represent the obstructions to passing from the local to the global.
By definition, if $0\to\sF\to\sG\to\sH\to0$ is an exact sequence of
abelian sheaves, then there is a long exact cohomology sequence:
\[\xymatrix{
  0 \ar[r]
    & \h^0(X,\sF) \ar[r]
    & \h^0(X,\sG) \ar[r]
    & \h^0(X,\sH) \ar[r]
    & \h^1(X,\sF) \ar[r]
    & \cdots \\
  \ldots \ar[r]
    & \h^n(X,\sF) \ar[r]
    & \h^n(X,\sG) \ar[r]
    & \h^n(X,\sH) \ar[r]
    & \h^{n+1}(X,\sF) \ar[r]
    & \cdots
}\]

\subsubsection{}\label{I:1-6-6}

Given an abelian sheaf $\sF$ on $\cS$, an \emph{$\sF$-torsor} is a
sheaf $\sG$ equipped with an action $\sF\times\sG\to\sG$ of $\sF$
such that locally (after restriction to all the objects of a covering
sieve of the final object $X$), $\sG$ with its $\sF$-action is
isomorphic to $\sF$ with the canonical translation action
$\sF\times\sF\to\sF$.

One can show that $\h^1(X,\sF)$ may be interpreted as the set of
isomorphism classes of $\sF$-torsors.

\section{The étale topology}\label{I:2}

We specialize the definitions of the preceding chapter to the étale
topology of a scheme $X$ (\ref{I:2-1}, \ref{I:2-2}, \ref{I:2-3}). When
$X$ is the spectrum of a field $K$, the corresponding cohomology agrees
with the Galois cohomology of $K$ (\ref{I:2-4}).

\subsection{The étale topology}\label{I:2-1}

We begin by recalling some facts about étale morphisms.

\begin{definition}\label{I:2-1-1}
Let $A$ be a (commutative) ring. An $A$-algebra $B$ is called étale if
$B$ is an $A$-algebra of finite presentation and the following equivalent
conditions hold:
\begin{enumerate}[\indent a)]
  \item For every $A$-algebra $C$ and every square-zero ideal $J$ of $C$,
    the canonical map
    \[
      \hom_{A\textnormal{-alg}}(B,C) \to \hom_{A\textnormal{-alg}}(B,C/J)
    \]
    is a bijection.
  \item $B$ is a flat $A$-module and $\Omega_{B/A}=0$ (where
    $\Omega_{B/A}$ denotes the module of relative differentials).
  \item Let $B=A[X_1,\dotsc,X_n]/I$ be a presentation of $B$. Then, for
    every prime ideal $\fp$ of $A[X_1,\dotsc,X_n]$ containing $I$, there
    exist polynomials $P_1,\dotsc,P_n\in I$ such that $I_\fp$ is generated
    by the images of $P_1,\dotsc,P_n$ and
    $\det(\partial P_i/\partial X_j)\notin\fp$.
\end{enumerate}
\end{definition}
(cf. \cite[I]{sga1} or \cite[V]{ra70})

A morphism of schemes $f:X\to S$ is called \emph{étale} if, for every
$x\in X$, there exist an affine open neighbourhood $U=\spec(A)$ of
$f(x)$ and an affine open neighbourhood $V=\spec(B)$ of $x$ in
$X\times_SU$ such that $B$ is an étale $A$-algebra.

\subsubsection{Examples}\label{I:2-1-2}
\begin{enumerate}[\indent a)]
  \item If $A$ is a field, an $A$-algebra $B$ is étale if and only if it
    is a finite product of separable extensions of $A$.
  \item If $X$ and $S$ are schemes of finite type over $\dC$, a morphism
    $f:X\to S$ is étale if and only if its analytification
    $f^{\text{an}}:X^{\text{an}}\to S^{\text{an}}$ is a local
    isomorphism.
\end{enumerate}

\subsubsection{Consequences}\label{I:2-1-3}
\begin{enumerate}[\indent a)]
  \item (base change) If $f:X\to S$ is an étale morphism, then so is
    $f_{S'}:X\times_SS'\to S'$ for every morphism $S'\to S$.
  \item If $f:X\to S$ and $g:Y\to S$ are two étale morphisms, then every
    $S$-morphism from $X$ to $Y$ is étale.
  \item (descent) Let $f:X\to S$ be a morphism. If there exists a
    faithfully flat morphism $S'\to S$ such that
    $f_{S'}:X\times_SS'\to S'$ is étale, then $f$ is étale.
\end{enumerate}

\subsubsection{}\label{I:2-1-4}

Let $X$ be a scheme, and let $\cS$ be the category of étale
$X$-schemes. By (\ref{I:2-1-3}.c), every morphism in $\cS$ is étale.
The \emph{étale topology} on $\cS$ is the topology for which a sieve on
$U$ is covering if it is generated by a finite family of morphisms
$\varphi_i:U_i\to U$ such that the union of the images of the
$\varphi_i$ covers $U$. The \emph{étale site} of $X$, denoted by
$\et X$, is the site defined by $\cS$ with the étale topology.

\subsection{Examples of sheaves}\label{I:2-2}

\subsubsection{Constant sheaf}\label{I:2-2-1}

Let $C$ be an abelian group and, for simplicity, suppose that $X$ is
noetherian. We denote by $\const C_X$ (or simply $C$ if there is no
ambiguity) the sheaf defined by $U\mapsto C^{\pi_0(U)}$, where
$\pi_0(U)$ is the (finite) set of connected components of $U$. The most
important case will be $C=\dZ/n$. Thus, by definition,
\[
  \h^0(X,\dZ/n)=\left(\dZ/n\right)^{\pi_0(X)}\text{.}
\]
Moreover, $\h^1(X,\dZ/n)$ is the set of isomorphism classes of
$\dZ/n$-torsors (\ref{I:1-6-6}), in other words, of finite Galois étale
coverings of $X$ with group $\dZ/n$. In particular, if $X$ is connected
and $\pi_1(X)$ is its fundamental group for a chosen base point, then
\[
  \h^1(X,\dZ/n) = \hom(\pi_1(X),\dZ/n)\text{.}
\]

\subsubsection{Multiplicative group}\label{I:2-2-2}

We denote by $\dG_{m,X}$ (or $\dG_m$ if there is no ambiguity) the sheaf
defined by $U\mapsto\Gamma(U,\sO_U^\times)$; it is indeed a sheaf by the
faithfully flat descent theorem (\ref{I:1-4-5}). By definition,
\[
  \h^0(X,\dG_m) = \h^0(X,\sO_X)^\times\text{;}
\]
in particular, if $X$ is reduced, connected, and proper over an
algebraically closed field $k$, then
\[
  \h^0(X,\dG_m) = k^\times\text{.}
\]

\begin{proposition}\label{I:2-2-3}
There is an isomorphism
\[
  \h^1(X,\dG_m) = \pic(X)\text{,}
\]
where $\pic(X)$ is the group of isomorphism classes of invertible sheaves
on $X$.
\end{proposition}
\begin{proof}
Let $*$ be the functor that associates to an invertible sheaf $\sL$ on
$X$ the following presheaf $\sL^*$ on $\et X$: for
$\varphi:U\to X$ étale,
\[
  \sL^*(U)=\operatorname{Isom}_U(\sO_U,\varphi^*\sL)\text{.}
\]
By (\ref{I:1-4-2}.i) and (\ref{I:1-4-5}) (full faithfulness), this
presheaf is a sheaf; indeed, it is a $\dG_m$-torsor. One checks
immediately that
\begin{enumerate}[\indent a)]
  \item the functor $*$ is compatible with (étale) localization;
  \item it induces an equivalence from the category of trivial invertible
    sheaves (i.e. isomorphic to $\sO_X$) to the category of trivial
    $\dG_m$-torsors: $\sL$ is trivial if and only if $\sL^*$ is trivial.
\end{enumerate}

Moreover, by (\ref{I:1-4-2}.ii) and (\ref{I:1-4-5}),
\begin{enumerate}[\indent a)]
\setcounter{enumi}{2}
  \item the notion of an invertible sheaf is local for the étale topology.
\end{enumerate}

It follows formally from a), b), and c) that $*$ is an equivalence between
the category of invertible sheaves on $X$ and that of $\dG_m$-torsors on
$\et X$; it induces the desired isomorphism. The inverse equivalence is
constructed as follows. If $\sT$ is a $\dG_m$-torsor, there exists a
finite étale covering $\{U_i\}$ of $X$ such that the torsors
$\sT/U_i$ are trivial; $\sT$ is then trivial on every $V$ étale over
$X$ belonging to the sieve $\cU\subset\et X$ generated by $\{U_i\}$.
On every $V\in\cU$, $\sT|V$ corresponds to an invertible sheaf $\sL_V$
(by b)), and the $\sL_V$ form an invertible sheaf given $\cU$-locally,
$\sL_\cU$ (by a)). By c), the latter comes from an invertible sheaf
$\sL(\sT)$ on $X$, and $\sT\mapsto\sL(\sT)$ is the desired inverse
of $*$.
\end{proof}

\subsubsection{Roots of unity}\label{I:2-2-4}

For every integer $n>0$, the sheaf of $n$th roots of unity, denoted by
$\dmu_n$, is the kernel of the $n$th-power map on $\dG_m$.
If $X$ is a scheme over a separably closed field $k$ and $n$ is
invertible in $k$, the choice of a primitive $n$th root of unity
$\zeta\in k$ defines an isomorphism $i\mapsto\zeta^i$ from $\dZ/n$ to
$\dmu_n$.

The relation between cohomology with coefficients in $\dmu_m$ and
cohomology with coefficients in $\dG_m$ is given by the cohomology exact
sequence deduced from the following theorem.

\begin{theorem}[Kummer theory]\label{I:2-2-5}
If $n$ is invertible on $X$, the $n$th-power map on $\dG_m$ is an
epimorphism of sheaves. Hence there is an exact sequence
\[
  0 \to \dmu_n \to \dG_m \to \dG_m \to 0\text{.}
\]
\end{theorem}
\begin{proof}
Let $U\to X$ be an étale morphism and let
$a\in\dG_m(U)=\Gamma(U,\sO_U^\times)$. Since $n$ is invertible on $U$,
the equation $T^n-a=0$ is separable; in other words,
$U'=\spec\left(\sO_U[T]/(T^n-a)\right)$ is étale over $U$. Moreover,
$U'\to U$ is surjective and this element has an $n$th root on $U'$, which proves
the result.
\end{proof}

\subsection{Fibres, direct images}\label{I:2-3}

\subsubsection{}\label{I:2-3-1}

A \emph{geometric point} of $X$ is a morphism $\bar x\to X$, where
$\bar x$ is the spectrum of a separably closed field $k(\bar x)$. By
abuse of notation, we denote it by $\bar x$, leaving the morphism
$\bar x\to X$ implicit. If $x$ is the image of $\bar x$ in $X$, we say
that $\bar x$ is centred at $x$. If $k(\bar x)$ is an algebraic extension
of the residue field $k(x)$, then $\bar x$ is called an
\emph{algebraic} geometric point of $X$.

An \emph{étale neighbourhood} of $\bar x$ is a commutative diagram
\[\xymatrix{
  & U \ar[d] \\
  \bar x \ar[ur] \ar[r]
  & X\text{,}
}\]
where $U\to X$ is étale.

The \emph{strict localization} of $X$ at $\bar x$ is the ring
$\sO_{X,\bar x}=\varinjlim\Gamma(U,\sO_U)$, where the direct limit is
taken over the étale neighbourhoods of $\bar x$. It is a strictly henselian
local ring whose residue field is the separable closure in $k(\bar x)$
of the residue field $k(x)$ of $X$ at $x$. It plays the role of the local
ring for the étale topology.

\subsubsection{}\label{I:2-3-2}

Given a sheaf $\sF$ on $\et X$, the \emph{fibre} of $\sF$ at $\bar x$
is the set (resp. group, \dots)
$\sF_{\bar x}=\varinjlim\sF(U)$, where the direct limit is again taken
over the étale neighbourhoods of $\bar x$.

A homomorphism of sheaves $\sF\to\sG$ is a mono-/epi-/isomorphism if and
only if the induced morphisms $\sF_{\bar x}\to\sG_{\bar x}$ on fibres
have the corresponding property at every geometric point of $X$. If $X$
is of finite type over an algebraically closed field, it suffices that
this hold at the rational points of $X$.

\subsubsection{}\label{I:2-3-3}

If $f:X\to Y$ is a morphism of schemes and $\sF$ is a sheaf on
$\et X$, the \emph{direct image} $f_*\sF$ of $\sF$ under $f$ is the
sheaf on $\et Y$ defined by
$f_*\sF(V)=\sF(X\times_YV)$ for every $V$ étale over $Y$. The functor
$f_*:\sh(\et X)\to\sh(\et Y)$ is left exact. Its right derived functors
$\eR^qf_*$ are called the higher direct images. If $\bar y$ is a
geometric point of $Y$, then
\[
  (\eR^q f_* \sF)_{\bar y} = \varinjlim \h^q(V\times_Y X,\sF)\text{,}
\]
where the direct limit is taken over the étale neighbourhoods $V$ of
$\bar y$.

Let $\sO_{Y,\bar y}$ be the strict localization of $Y$ at $\bar y$,
let $\widetilde Y=\spec(\sO_{Y,\bar y})$, and let
$\widetilde X=X\times_Y\widetilde Y$. The sheaf $\sF$ may be extended
to $\et{\widetilde X}$ (this is a special case of the general notion of
inverse image) as follows. Let $\widetilde U$ be a scheme étale over
$\widetilde X$. Then there exist an étale neighbourhood $V$ of $\bar y$
and a scheme $U$ étale over $X\times_YV$ such that
$\widetilde U=U\times_V\widetilde Y$; set
\[
  \sF(\widetilde U)=\varinjlim \sF\left(U\times_V V'\right)\text{,}
\]
where the direct limit is taken over the étale neighbourhoods $V'$ of
$\bar y$ that dominate $V$. With this definition,
\[
  \left(\eR^q f_* \sF\right)_{\bar y} = \h^q(\widetilde X,\sF)\text{.}
\]

The functor $f_*$ has a left adjoint $f^*$, the “inverse image” functor.
If $\bar x$ is a geometric point of $X$ and $f(\bar x)$ its image in
$Y$, then $(f^*\sF)_{\bar x}=\sF_{f(\bar x)}$. This formula shows that
$f^*$ is exact. Consequently, $f_*$ sends injective sheaves to injective
sheaves, and the spectral sequence for the composite functor
$\Gamma\circ f_*$ (resp. $g_*f_*$) gives the following theorem.

\begin{theorem}[Leray spectral sequence]\label{I:2-3-4}
Let $\sF$ be an abelian sheaf on $\et X$ and let $f:X\to Y$ be a
morphism of schemes (resp. let
$X\xrightarrow fY\xrightarrow gZ$ be morphisms of schemes). There is a
spectral sequence
\begin{align*}
  E_2^{pq} &= \h^p(Y,\eR^q f_* \sF) \Rightarrow \h^{p+q}(X,\sF) \\
  \text{(resp.}\qquad E_2^{pq} &= \eR^p g_* \eR^q f_* \sF \Rightarrow \eR^{p+q}(gf)_* \sF\text{).}
\end{align*}
\end{theorem}

\begin{corollary}\label{I:2-3-5}
If $\eR^qf_*\sF=0$ for every $q>0$, then
$\h^p(Y,f_*\sF)=\h^p(X,\sF)$ (resp.
$\eR^pg_*(f_*\sF)=\eR^p(gf)_*\sF$) for every $p\geqslant0$.
\end{corollary}

This applies in particular in the following case:

\begin{proposition}\label{I:2-3-6}
Let $f:X\to Y$ be a finite morphism (or even, by passage to the limit,
an integral morphism), and let $\sF$ be an abelian sheaf on $X$. Then
$\eR^qf_*\sF=0$ for every $q>0$.
\end{proposition}

Indeed, let $\bar y$ be a geometric point of $Y$, let $\widetilde Y$ be
the spectrum of the strict localization of $Y$ at $\bar y$, and let
$\widetilde X=X\times_Y\widetilde Y$. By the preceding discussion, it
suffices to prove that $\h^q(\widetilde X,\sF)=0$ for every $q>0$. But
$\widetilde X$ is the spectrum of a product of strictly henselian local
rings (cf. \cite[I]{ra70}), and the functor
$\Gamma(\widetilde X,-)$ is exact because every surjective étale
$\widetilde X$-scheme admits a section. This proves the assertion.

\subsection{Galois cohomology}\label{I:2-4}

For $X=\spec(K)$, the spectrum of a field, we shall see that étale
cohomology can be identified with Galois cohomology.

\subsubsection{}\label{I:2-4-1}

We begin with a topological analogy. If $K$ is the function field of an
integral affine algebraic variety $Y=\spec(A)$ over $\dC$, then
$K=\varinjlim_{f\in A}A[1/f]$.

In other words, $X=\varprojlim U$, where $U$ runs over the open subsets
of $Y$. It is known that there exist arbitrarily small Zariski open
subsets that are $K(\pi,1)$ spaces for the classical topology. It is
therefore natural to regard $\spec(K)$ itself as a $K(\pi,1)$, where
$\pi$ is the fundamental group (in the algebraic sense) of $X$, namely
the Galois group of $\bar K/K$, where $\bar K$ is a separable closure of
$K$.

\subsubsection{}\label{I:2-4-2}

More precisely, let $K$ be a field, let $\bar K$ be a separable closure
of $K$, and let $G=\gal(\bar K/K)$ be the topological Galois group. To
every finite étale $K$-algebra $A$ (a finite product of separable
extensions of $K$), associate the finite set $\hom_K(A,\bar K)$. The
Galois group $G$ acts on this set through a discrete (hence finite)
quotient. If $A=K[T]/(F)$, this set is identified with the set of roots
of $F$ in $\bar K$. Galois theory, in the form given to it by
Grothendieck, says:

\begin{proposition}\label{I:2-4-3}
The functor
\[
  \left\{\begin{array}{c}
          \textnormal{$K$-algebras} \\
          \textnormal{finite étale}
        \end{array}\right\}
  \to
  \left\{\begin{array}{c}
          \textnormal{finite sets on which} \\
          \textnormal{$G$ acts continuously}
        \end{array}\right\}
\]
that associates $\hom_K(A,\bar K)$ to an étale algebra $A$ is an
anti-equivalence of categories.
\end{proposition}

This yields an analogous description of sheaves for the étale topology on
$\spec(K)$:

\begin{proposition}\label{I:2-4-4}
The functor
\[
  \left\{\begin{array}{c}
          \textnormal{étale sheaves} \\
          \textnormal{on $\spec(K)$}
        \end{array}\right\}
  \to
  \left\{\begin{array}{c}
          \textnormal{sets on which} \\
          \textnormal{$G$ acts continuously}
        \end{array}\right\}
\]
that associates to a sheaf $\sF$ its fibre $\sF_{\bar K}$ at the
geometric point $\spec(\bar K)$ is an equivalence of categories.
\end{proposition}

The group $G$ is said to act continuously on a set $E$ if the stabilizer
of every element of $E$ is an open subgroup of $G$. The inverse functor
is described in the evident way: let $A$ be a finite étale
$K$-algebra, let $U=\spec(A)$, and let
$U(\bar K)=\hom_K(A,\bar K)$ be the corresponding $G$-set; then
$\sF(U)=\hom_{G\text{-set}}(U(\bar K),\sF_{\bar K})$.

In particular, if $X=\spec(K)$, then $\sF(X)=\sF_{\bar K}^G$. Restricting
to abelian sheaves and passing to derived functors gives canonical
isomorphisms
\[
  \h^q(\et X,\sF) = \h^q(G,\sF_{\bar K})
\]

\subsubsection{Examples}\label{I:2-4-5}

\begin{enumerate}[\indent a)]
  \item The constant sheaf $\dZ/n$ corresponds to $\dZ/n$ with the
    trivial action of $G$.
  \item The sheaf $\dmu_n$ of $n$th roots of unity corresponds to the
    group $\dmu_n(\bar K)$ of $n$th roots of unity in $\bar K$, with
    the natural action of $G$.
  \item The sheaf $\dG_m$ corresponds to the group $\bar K^\times$ with
    the natural action of $G$.
\end{enumerate}

\section{The cohomology of curves}\label{I:3}

For topological spaces, unwindings using the K\"unneth formula and
simplicial decompositions reduce the calculation of cohomology to the
interval $I=[0,1]$, for which $\h^0(I,\dZ)=\dZ$ and
$\h^q(I,\dZ)=0$ for $q>0$.

In our case, the unwindings lead to more complicated objects, namely curves
over an algebraically closed field; in this chapter we calculate their
cohomology. The situation is more complicated than in the topological case,
because the cohomology groups vanish only for $q>2$. The essential
ingredient in the calculations is the vanishing of the Brauer group of the
function field of such a curve (Tsen's theorem, \ref{I:3-2}).

\subsection{The Brauer group}\label{I:3-1}

We first recall its classical definition.

\begin{definition}\label{I:3-1-1}
Let $K$ be a field and let $A$ be a finite-dimensional $K$-algebra. The
algebra $A$ is called a central simple algebra over $K$ if the following
equivalent conditions hold:
\begin{enumerate}[\indent a)]
  \item $A$ has no nontrivial two-sided ideal, and its centre is $K$.
  \item There exists a finite Galois extension $K'/K$ such that
    $A_{K'}=A\otimes_KK'$ is isomorphic to an algebra of square matrices
    over $K'$.
  \item $A$ is $K$-isomorphic to an algebra of square matrices over a
    division algebra with centre $K$.
\end{enumerate}
\end{definition}

Two such algebras are called equivalent if the division algebras associated
with them by c) are $K$-isomorphic. If the algebras have the same
dimension, this is equivalent to their being $K$-isomorphic. Tensor product
induces, on passing to the quotient, an abelian group structure on the set
of equivalence classes. This group is classically called the
\emph{Brauer group} of $K$ and is denoted by $\br(K)$.

\subsubsection{}\label{I:3-1-2}

We denote by $\br(n,K)$ the set of $K$-isomorphism classes of
$K$-algebras $A$ such that there exists a finite Galois extension $K'$
of $K$ for which $A_{K'}$ is isomorphic to the algebra $M_n(K')$ of
$n\times n$ square matrices over $K'$. By definition, $\br(K)$ is the
union of the subsets $\br(n,K)$ for $n\in\dN$. Let $\bar K$ be an
algebraic closure of $K$ and let $G=\gal(\bar K/K)$. The set $\br(n,K)$
is the set of “forms” of $M_n(\bar K)$, and hence is canonically
isomorphic to $\h^1\left(G,\aut(M_n(\bar K))\right)$.

Every automorphism of $M_n(\bar K)$ is inner. Consequently,
$\aut(M_n(\bar K))$ is identified with the projective linear group
$\pgl(n,\bar K)$, and there is a canonical bijection
\[
  \theta_n : \br(n,K) \iso \h^1\left(G,\pgl(n,\bar K)\right)\text{.}
\]

On the other hand, the exact sequence
\begin{equation*}\tag{3.1.2.1}\label{I:eq:3-1-2-1}
  1 \to \bar K^\times \to \gl(n,\bar K) \to \pgl(n,\bar K) \to 1\text{,}
\end{equation*}
defines a coboundary operator
\[
  \Delta_n : \h^1\left(G,\pgl(n,\bar K)\right) \to \h^2(G,\bar K^\times)\text{.}
\]

Composing $\theta_n$ and $\Delta_n$ gives a map
\[
  \delta_n : \br(n,K) \to \h^2(G,\bar K^\times)\text{.}
\]
It is readily checked that the maps $\delta_n$ are compatible and define
a group homomorphism
\[
  \delta : \br(K) \to \h^2(G,\bar K^\times)\text{.}
\]

\begin{proposition}\label{I:3-1-3}
The homomorphism $\delta:\br(K)\to\h^2(G,\bar K^\times)$ is bijective.
\end{proposition}

This follows from the next two lemmas.

\begin{lemma}\label{I:3-1-4}
The map
$\Delta_n:\h^1\left(G,\pgl(n,\bar K)\right)\to
\h^2(G,\bar K^\times)$ is injective.
\end{lemma}

By \cite{se94}, Corollary to Proposition I-44, it suffices to check that
whenever the exact sequence (\ref{I:eq:3-1-2-1}) is twisted by an element
of $\h^1\left(G,\pgl(n,\bar K)\right)$, the $\h^1$ of the middle group
is trivial. This middle group is the group of $\bar K$-points of the
multiplicative group of a central simple algebra $A$ of rank $n^2$ over
$K$. To prove that $\h^1(G,A_{\bar K}^\times)=0$, interpret $A^\times$
as the automorphism group of the free $A$-module $L$ of rank $1$, and
$\h^1$ as the set of “forms” of $L$---$A$-modules of rank $n^2$ over
$K$, which are automatically free.

\begin{lemma}\label{I:3-1-5}
Let $\alpha\in\h^2(G,\bar K^\times)$, let $K'$ be a finite extension of
$K$ contained in $\bar K$, let $n=[K':K]$, and let
$G'=\gal(\bar K/K')$. If the image of $\alpha$ in
$\h^2(G',\bar K^\times)$ is zero, then $\alpha$ belongs to the image of
$\Delta_n$.
\end{lemma}

First note that
\[
  \h^2(G',\bar K^\times) \simeq \h^2\left(G,(\bar K\otimes_K K')^\times\right)\text{.}
\]
(Geometrically, if $x=\spec(K)$, $x'=\spec(K')$, and
$\pi:x'\to x$ is the canonical morphism, then
$\eR^q\pi_*(\dG_{m,x'})=0$ for $q>0$ and consequently
$\h^q(x',\dG_{m,x'})\simeq\h^q(x,\pi_*\dG_{m,x'})$ for
$q\geqslant0$.)

Moreover, the choice of a basis of $K'$ as a vector space over $K$
defines a homomorphism
\[
  (\bar K\otimes_K K')^\times \to \gl(n,\bar K)
\]
that sends an element $x$ to the endomorphism of $\bar K\otimes_KK'$
given by multiplication by $x$. There is then a commutative diagram with
exact rows:
\[\xymatrix{
  1 \ar[r]
    & \bar K^\times \ar[r] \ar@{=}[d]
    & (\bar K\otimes_K K')^\times \ar[r] \ar[d]
    & (\bar K\otimes_K K')^\times / \bar K^\times \ar[r] \ar[d]
    & 1 \\
  1 \ar[r]
    & \bar K^\times \ar[r]
    & \gl(n,\bar K) \ar[r]
    & \pgl(n,\bar K) \ar[r]
    & 1
}\]
The lemma follows from the commutative diagram obtained by passing to
cohomology:
\[\xymatrix{
  \h^1\left(G,(\bar K\otimes_K K')^\times/\bar K^\times\right) \ar[r] \ar[d]
    & \h^2(G,\bar K^\times) \ar[r] \ar@{=}[d]
    & \h^2\left(G,(\bar K\otimes_K K')^\times\right) \\
  \h^1\left(G,\pgl(n,\bar K)\right) \ar[r]^-{\Delta_n}
    & \h^2(G,\bar K^\times) \text{.}
}\]

\begin{proposition}\label{I:3-1-6}
Let $K$ be a field, let $\bar K$ be an algebraic closure of $K$, and let
$G=\gal(\bar K/K)$. Suppose that $\br(K')=0$ for every finite extension
$K'$ of $K$. Then:
\begin{enumerate}[\indent i)]
  \item $\h^q(G,\bar K^\times)=0$ for every $q>0$.
  \item $\h^q(G,\sF)=0$ for every torsion $G$-module $\sF$ and every
    $q\geqslant2$.
\end{enumerate}
\end{proposition}

(For the proof, cf. \cite{se94}.)

\subsection{Tsen's theorem}\label{I:3-2}

\begin{definition}\label{I:3-2-1}
A field $K$ is called $C_1$ if every nonconstant homogeneous polynomial
$f(x_1,\dotsc,x_n)$ of degree $d<n$ has a nontrivial zero.
\end{definition}

\begin{proposition}\label{I:3-2-2}
If a field $K$ is $C_1$, then $\br(K)=0$.
\end{proposition}

We must show that every division algebra $D$ with centre $K$ and finite
over $K$ is equal to $K$. Let $r^2$ be the degree of $D$ over $K$ and
let $\operatorname{Nrd}:D\to K$ be the reduced norm.

(Locally for the étale topology on $K$, $D$ is isomorphic---not
canonically---to a matrix algebra $M_r$, and the reduced norm agrees with
the determinant map. This is well defined independently of the chosen
isomorphism between $D$ and $M_r$, because every automorphism of $M_r$
is inner and similar matrices have the same determinant. This map, defined
locally for the étale topology, descends, by its local uniqueness, to a map
$\operatorname{Nrd}:D\to K$.)

The only zero of $\operatorname{Nrd}$ is the zero element of $D$, since
for $x\ne0$ one has
$\operatorname{Nrd}(x)\cdot\operatorname{Nrd}(x^{-1})=1$. On the other
hand, if $\{e_1,\dotsc,e_{r^2}\}$ is a basis of $D$ over $K$ and
$x=\sum x_ie_i$, then $\operatorname{Nrd}(x)$ is a homogeneous
polynomial $\operatorname{Nrd}(x_1,\dotsc,x_{r^2})$ of degree $r$
(this is clear locally for the étale topology). Since $K$ is $C_1$, one
has $r^2\leqslant r$, that is, $r=1$ and $D=K$.

\begin{theorem}[Tsen]\label{I:3-2-3}
Let $k$ be an algebraically closed field and let $K$ be an extension of
$k$ of transcendence degree $1$. Then $K$ is $C_1$.
\end{theorem}

First suppose that $K=k(X)$. Let
\[
  f(T) = \sum a_{i_1,\dotsc i_n} T_1^{i_1} \dotsm T_n^{i_n}
\]
be a homogeneous polynomial of degree $d<n$ with coefficients in
$k(X)$. After multiplying the coefficients by a common denominator, we
may suppose that they lie in $k[X]$. Let
$\delta=\sup\deg(a_{i_1\dotsc i_n})$. We seek a nontrivial zero in
$k[X]$ by the method of undetermined coefficients, writing each $T_i$
($i=1,\dotsc,n$) as a polynomial of degree $N$ in $X$. The equation
$f(T)=0$ then becomes a system of homogeneous equations in the
$n\times(N+1)$ coefficients of the polynomials $T_i(X)$, expressing the
vanishing of the coefficients of the polynomial in $X$ obtained by
replacing $T_i$ with $T_i(X)$. This polynomial has degree at most
$\delta+Nd$, so there are $\delta+Nd+1$ equations in
$n\times(N+1)$ variables. Since $k$ is algebraically closed, this system
has a nontrivial solution if $n(N+1)>Nd+\delta+1$, which holds for
sufficiently large $N$ when $d<n$.

To prove the theorem in general, it plainly suffices to prove it when
$K$ is a finite extension of the purely transcendental extension
$k(X)$ of $k$. Let $f(T)=f(T_1,\dotsc,T_n)$ be a homogeneous
polynomial of degree $d<n$ with coefficients in $K$. Let
$s=[K:k(X)]$ and let $e_1,\dotsc,e_s$ be a basis of $K$ over $k(X)$.
Introduce $sn$ new variables $U_{ij}$ such that
$T_i=\sum U_{ij}e_j$. For $f(T)$ to have a nontrivial zero in $K$, it
suffices that the polynomial $g(U_{ij})=N_{K/k(X)}(f(T))$ have a nontrivial
zero in $k(X)$. But $g$ is a polynomial of degree $sd$ in $sn$
variables, which proves the result.

\begin{corollary}\label{I:3-2-4}
Let $k$ be an algebraically closed field and let $K$ be an extension of
$k$ of transcendence degree $1$. Then the étale cohomology groups
$\h^q(\spec(K),\dG_m)$ vanish for every $q>0$.
\end{corollary}

\subsection{Cohomology of smooth curves}

From now on, unless expressly stated otherwise, all cohomology groups are
étale cohomology groups.

\begin{proposition}\label{I:3-3-1}
Let $k$ be an algebraically closed field and let $X$ be a connected
nonsingular projective curve over $k$. Then
\begin{align*}
  \h^0(X,\dG_m) &= k^\times \text{,} \\
  \h^1(X,\dG_m) &= \pic(X) \text{,} \\
  \h^q(X,\dG_m) &= 0 \text{ for $q\geqslant 2$.}
\end{align*}
\end{proposition}

Let $\eta$ be the generic point of $X$, let $j:\eta\to X$ be the
canonical morphism, and let $\dG_{m,\eta}$ be the multiplicative group of
the function field $K(X)$. For every closed point $x$ of $X$, let
$i_x:x\to X$ be the canonical immersion and let $\dZ_x$ be the constant
sheaf with value $\dZ$ on $x$. Thus $j_*\dG_{m,\eta}$ is the sheaf of
nonzero meromorphic functions on $X$ and
$\oplus_{x\in X}i_{x*}\dZ_x$ is the sheaf of divisors; hence there is an
exact sequence of sheaves:
\begin{equation*}\tag{3.3.1.1}\label{I:eq:3-3-1-1}
\xymatrix{
  0 \ar[r]
    & \dG_m \ar[r]
    & j_* \dG_{m,\eta} \ar[r]^-{\dv}
    & \displaystyle\bigoplus_{x\in X} i_{x*} \dZ_x \ar[r]
    & 0\text{.}
}
\end{equation*}

\begin{lemma}\label{I:3-3-2}
One has $\eR^qj_*\dG_{m,\eta}=0$ for every $q>0$.
\end{lemma}

It suffices to show that the fibre of this sheaf at every closed point
$x$ of $X$ is zero. If $\tilde\sO_{X,x}$ is the henselization of $X$ at
$x$ and $K$ is the fraction field of $\tilde\sO_{X,x}$, then
\[
  \spec(K) = \eta\times_X \spec(\tilde\sO_{X,x})\text{,}
\]
and hence $(\eR^qj_*\dG_{m,\eta})_x=\h^q(\spec(K),\dG_m)$.

Now $K$ is an algebraic extension of $k(X)$, hence an extension of $k$
of transcendence degree $1$; the lemma follows from (\ref{I:3-2-4}).

\begin{lemma}\label{I:3-3-3}
One has $\h^q(X,j_*\dG_{m,\eta})=0$ for every $q>0$.
\end{lemma}

Indeed, (\ref{I:3-3-2}) and the Leray spectral sequence for $j$ give
\[
  \h^q(X, j_* \dG_{m,\eta}) = \h^q(\eta,\dG_{m,\eta})
\]
for every $q\geqslant0$, and the right-hand side vanishes for $q>0$ by
(\ref{I:3-2-4}).

\begin{lemma}\label{I:3-3-4}
One has
$\h^q\left(X,\bigoplus_{x\in X}i_{x*}\dZ_x\right)=0$ for every $q>0$.
\end{lemma}

Indeed, for every closed point $x$ of $X$ one has
$\eR^qi_{x*}\dZ_x=0$ for $q>0$, because $i_x$ is a finite morphism
(\ref{I:2-3-6}), and
\[
  \h^q(X,i_{x*}\dZ_x) = \h^q(x,\dZ_x)\text{.}
\]
The right-hand side vanishes for every $q>0$ because $x$ is the spectrum
of an algebraically closed field. (This also shows that the lemma holds
more generally for every skyscraper sheaf on $X$.)

The preceding lemmas and the exact sequence (\ref{I:eq:3-3-1-1}) give
\[
  \h^q(X,\dG_m) = 0 \text{ for $q\geqslant 2$,}
\]
and the low-degree exact cohomology sequence
\[
  1 \to \h^0(X,\dG_m) \to \h^0(X, j_*\dG_{m,\eta}) \to \h^0\left(X,\bigoplus_{x\in X} i_{x*}\dZ_x\right) \to \h^1(X,\dG_m) \to 1
\]
which is precisely the exact sequence
\[
  1 \to k^\times \to k(X)^\times\to \Div(X) \to \pic(X) \to 1\text{.}
\]
Proposition \ref{I:3-3-1} implies that the cohomology groups of $X$ with
values in $\dZ/n$, for $n$ prime to the characteristic of $k$, have
reasonable values:

\begin{corollary}\label{I:3-3-5}
If $X$ has genus $g$ and $n$ is invertible in $k$, then
$\h^q(X,\dZ/n)$ vanishes for $q>2$ and is free over $\dZ/n$ of ranks
$1,2g,1$ for $q=0,1,2$. Replacing $\dZ/n$ by the isomorphic group
$\dmu_n$, there are canonical isomorphisms
\begin{align*}
  \h^0(X,\dmu_n) &= \dmu_n \\
  \h^1(X,\dmu_n) &= \pic^0(X)_n \\
  \h^2(X,\dmu_n) &= \dZ/n \text{.}
\end{align*}
\end{corollary}

Since $k$ is algebraically closed, $\dZ/n$ is isomorphic
(noncanonically) to $\dmu_n$. The Kummer exact sequence
\[
  0 \to \dmu_n \to \dG_m \to \dG_m \to 0\text{,}
\]
and Proposition \ref{I:3-3-1} yield the equalities
\[
  \h^q(X,\dZ/n) = 0 \text{ for $q>2$,}
\]
and, in low degrees, the exact sequences
\[\xymatrix{
  0 \ar[r]
    & \h^0(X,\dmu_n) \ar[r]
    & k^\times \ar[r]^-n
    & k^\times \ar[r]
    & 0
}\]
\[\xymatrix{
  0 \ar[r]
    & \h^1(X,\dmu_n) \ar[r]
    & \pic(X) \ar[r]^-n
    & \h^2(X,\dmu_n) \ar[r]
    & 0\text{.}
}\]
Moreover, there is an exact sequence
\[\xymatrix{
  0 \ar[r]
    & \pic^0(X) \ar[r]
    & \pic(X) \ar[r]^-{\deg}
    & \dZ \ar[r]
    & 0 \text{,}
}\]
and $\pic^0(X)$ is identified with the group of $k$-rational points of
an abelian variety of dimension $g$, the Jacobian of $X$. In such a
group, multiplication by $n$ is surjective and its kernel is a free
$\dZ/n\dZ$-module of rank $2g$ (because $n$ is invertible in $k$);
this proves the corollary.

A clever unwinding using the “trace method” yields the following
proposition as a corollary.

\begin{proposition}[{\cite[IX 5.7]{sga4}}]\label{I:3-3-6}
Let $k$ be an algebraically closed field, let $X$ be an algebraic curve
over $k$, and let $\sF$ be a torsion sheaf on $X$. Then:
\begin{enumerate}[\indent i)]
  \item $\h^q(X,\sF)=0$ for $q>2$.
  \item If $X$ is affine, then even $\h^q(X,\sF)=0$ for $q>1$.
\end{enumerate}
\end{proposition}

For the proof, as well as for an exposition of the “trace method,” we
refer to \cite[IX 5]{sga4}.

\subsection{Unwindings}\label{I:3-4}

To calculate the cohomology of varieties of dimension $>1$, one uses
fibrations by curves, thereby reducing to the study of morphisms whose
fibres have dimension $\leqslant1$. This principle has several variants;
we mention some of them.

\subsubsection{}\label{I:3-4-1}

Let $A$ be a $k$-algebra of finite type, and let $a_1,\dotsc,a_n$ be
generators of $A$. If we put $X_0=\spec(k)$,
$X_i=\spec(k[a_1,\dotsc,a_i])$, and $X_n=\spec(A)$, the canonical
inclusions
$k[a_1,\dotsc,a_i]\to k[a_1,\dotsc,a_i,a_{i+1}]$ define morphisms
$X_n\to X_{n-1}\to\cdots\to X_1\to X_0$ whose fibres have dimension
$\leqslant1$.

\subsubsection{}\label{I:3-4-2}

For a smooth morphism, one can be more precise. An
\emph{elementary fibration} is a morphism of schemes $f:X\to S$ that
can be embedded in a commutative diagram
\[\xymatrix{
  X \ar@{^{(}->}[r]^-j \ar[dr]_-f
    & \bar X \ar[d]^-{\bar f}
    & \ar@{_{(}->}[l]_-i \ar[dl]^-g Y \\
  & S
}\]
satisfying the following conditions:
\begin{enumerate}[\indent i)]
  \item $j$ is an open immersion, dense in every fibre, and
    $X=\bar X\setminus Y$.
  \item $\bar f$ is smooth and projective, with irreducible geometric
    fibres of dimension $1$.
  \item $g$ is an étale covering and no fibre of $g$ is empty.
\end{enumerate}

A \emph{good neighbourhood} relative to $S$ is an $S$-scheme $X$ for
which there exist $S$-schemes $X=X_n,\dotsc,X_0=S$ and elementary
fibrations $f_i:X_i\to X_{i-1}$, $i=1,\dotsc,n$. One can show
\cite[XI 3.3]{sga4} that if $X$ is a smooth scheme over an algebraically
closed field $k$, then every rational point of $X$ has an open
neighbourhood that is a good neighbourhood (relative to $\spec(k)$).

\subsubsection{}\label{I:3-4-3}

A proper morphism $f:X\to S$ can be unwound as follows. By Chow's lemma,
there is a commutative diagram
\[\xymatrix{
  X \ar[dr]_-f
    & \ar[l]_-\pi \ar[d]^-{\bar f} \bar X \\
  & S
}\]
in which $\pi$ and $\bar f$ are projective morphisms, with $\pi$ in
addition an isomorphism over a dense open subset of $X$. Locally on
$S$, $\bar X$ is a closed subscheme of a standard projective space
$\dP_S^n$.

The latter is unwound by considering the projection
$\varphi:\dP_S^n\to\dP_S^1$ that sends the point with homogeneous
coordinates $(x_0,x_1,\dotsc,x_n)$ to $(x_0,x_1)$. This is a rational
map defined outside the closed subset $Y\simeq\dP_S^{n-2}$ of
$\dP_S^n$ given by the homogeneous equations $x_0=x_1=0$. Let
$u:P\to\dP_S^n$ be the blow-up with centre $Y$; the fibres of $u$ have
dimension $\leqslant1$. Moreover, there is a natural morphism
$v:P\to\dP_S^1$ extending the rational map $\varphi$, and $v$ makes
$P$ a $\dP_S^1$-scheme locally isomorphic to the standard projective
space $\dP^{n-1}$, which may in turn be projected onto a $\dP^1$, and
so on.

\subsubsection{}\label{I:3-4-4}

A smooth projective variety $X$ can be swept out by a
\emph{Lefschetz pencil}. The blow-up $\widetilde X$ of the intersection
of the axis of the pencil with $X$ maps to $\dP^1$, and the fibres of
this map are the hyperplane sections of $X$ by the hyperplanes of the
pencil.

\section{The proper base change theorem}\label{I:4}

\subsection{Introduction}\label{I:4-1}

This chapter is devoted to the proof and applications of the following
theorem.

\begin{theorem}\label{I:4-1-1}
Let $f:X\to S$ be a proper morphism of schemes and let $\sF$ be a
torsion abelian sheaf on $X$. Then, for every $q\geqslant0$, the fibre
of $\eR^qf_*\sF$ at a geometric point $s$ of $S$ is isomorphic to the
cohomology $\h^q(X_s,\sF)$ of the fibre
$X_s=X\otimes_S\spec k(s)$ of $f$ at $s$.
\end{theorem}

For a proper separated continuous map $f:X\to S$ between topological
spaces (separated means that the diagonal of $X\times_SX$ is closed) and
an abelian sheaf $\sF$ on $X$, the analogous result is well known and
elementary. Since $f$ is closed, the $f^{-1}(V)$, with $V$ ranging over
the neighbourhoods of $s$, form a fundamental system of neighbourhoods
of $X_s$, and one checks that
$\h^\bullet(X_s,\sF)=\varinjlim_U\h^\bullet(U,\sF)$, where $U$ ranges
over the neighbourhoods of $X_s$. In practice, $X_s$ even has a
fundamental system $\cU$ of neighbourhoods $U$ of which it is a
deformation retract, and for $\sF$ constant one therefore has
$\h^\bullet(X_s,\sF)=\h^\bullet(U,\sF)$. Figuratively speaking, the
special fibre swallows the general fibre.

For schemes the proof is more delicate, and it is essential to suppose
that $\sF$ is torsion (\cite[XII.2]{sga4}). In view of the description of
the fibres of $\eR^qf_*\sF$ (\ref{I:2-3-3}), Theorem
(\ref{I:4-1-1}) is essentially equivalent to the following theorem.

\begin{theorem}\label{I:4-1-2}
Let $A$ be a strictly henselian local ring and let $S=\spec(A)$. Let
$f:X\to S$ be proper, and let $X_0$ be the closed fibre of $f$. Then, for every
torsion abelian sheaf $\sF$ on $X$ and every $q\geqslant0$, one has
$\h^q(X,\sF)\iso\h^q(X_0,\sF)$.
\end{theorem}

By passage to the limit, it suffices to prove the theorem when $A$ is the
strict henselization of a finitely generated $\dZ$-algebra at a prime
ideal. We first treat the case $q=0$ or $1$ and $\sF=\dZ/n$
(\ref{I:4-2}). An argument based on constructible sheaves (\ref{I:4-3})
shows moreover that it suffices to consider constant $\sF$. On the other
hand, the unwinding in (\ref{I:3-4-3}) allows us to suppose that $X_0$
is a curve; in that case it remains only to prove the theorem for $q=2$
(\ref{I:4-4}).

Among other applications (\ref{I:4-6}), the theorem allows one to define
cohomology with proper support (\ref{I:4-5}).

\subsection{Proof for \texorpdfstring{$q=0$}{q=0} or \texorpdfstring{$1$}{1} and \texorpdfstring{$\sF=\dZ/n$}{\sF=Z/n}}\label{I:4-2}

The result for $q=0$ and $\sF$ constant is equivalent to the following
proposition (Zariski's connectedness theorem).

\begin{proposition}\label{I:4-2-1}
Let $A$ be a noetherian henselian local ring and let $S=\spec(A)$. Let
$f:X\to S$ be a proper morphism, and let $X_0$ be the closed fibre of
$f$. Then the sets of connected components $\pi_0(X)$ and
$\pi_0(X_0)$ are in bijection.
\end{proposition}

Equivalently, we must show that the sets $\operatorname{Of}(X)$ and
$\operatorname{Of}(X_0)$ of open-and-closed subsets are in bijection.
The set $\operatorname{Of}(X)$ is known to correspond bijectively to the
set of idempotents of $\Gamma(X,\sO_X)$, and similarly
$\operatorname{Of}(X_0)$ corresponds bijectively to the set of
idempotents of $\Gamma(X_0,\sO_{X_0})$. Thus we must show that the
canonical map
\[
  \operatorname{Idem} \Gamma(X,\sO_X) \to \operatorname{Idem}\Gamma(X_0,\sO_{X_0})
\]
is bijective.

Let $\fm$ denote the maximal ideal of $A$, let
$\Gamma(X,\sO_X)^\wedge$ denote the completion of $\Gamma(X,\sO_X)$
for the $\fm$-adic topology, and, for every integer $n\geqslant0$, let
$X_n=X\otimes_AA/\fm^{n+1}$. By the finiteness theorem for proper
morphisms \cite[III.3.2]{ega}, $\Gamma(X,\sO_X)$ is a finite
$A$-algebra; since $A$ is henselian, the canonical map
\[
  \operatorname{Idem} \Gamma(X,\sO_X) \to \operatorname{Idem}\Gamma(X,\sO_X)^\wedge
\]
is bijective. In particular, the canonical map
\[
  \operatorname{Idem} \Gamma(X,\sO_X)^\wedge \to \varprojlim \operatorname{Idem}\Gamma(X_n,\sO_{X_n})
\]
is bijective. But since $X_n$ and $X_0$ have the same underlying
topological space, the canonical map
\[
  \operatorname{Idem}\Gamma(X_n,\sO_{X_n}) \to \operatorname{Idem} \Gamma(X_0,\sO_{X_0})
\]
is bijective for every $n$, which completes the proof.

Since $\h^1(X,\dZ/n)$ is in bijection with the set of isomorphism classes
of Galois étale coverings of $X$ with group $\dZ/n$, the theorem for
$q=1$ and $\sF=\dZ/n$ follows from the next proposition.

\begin{proposition}\label{I:4-2-2}
Let $A$ be a noetherian henselian local ring and let $S=\spec(A)$. Let
$f:X\to S$ be a proper morphism, and let $X_0$ be the closed fibre of $f$. Then
the restriction functor
\[
  \operatorname{Rev.et}(X) \to \operatorname{Rev.et}(X_0)
\]
is an equivalence of categories.
\end{proposition}

(If $X_0$ is connected and a geometric point of $X_0$ has been chosen
as base point, this says that the canonical map
$\pi_1(X_0)\to\pi_1(X)$ on (profinite) fundamental groups is bijective.)

Proposition (\ref{I:4-2-1}) shows that this functor is fully faithful.
Indeed, if $X'$ and $X''$ are two étale coverings of $X$, an
$X$-morphism from $X'$ to $X''$ is determined by its graph, which is an
open-and-closed subset of $X'\times_XX''$.

It remains to show that every étale covering $X_0'$ of $X_0$ extends to
an étale covering of $X$. Étale coverings are insensitive to nilpotent
elements \cite[ch.1]{sga1}; consequently, $X_0'$ lifts uniquely to an
étale covering $X_n'$ of $X_n$ for every $n\geqslant0$, in other words,
to an étale covering $\fX'$ of the formal scheme $\fX$ obtained by
completing $X$ along $X_0$. By Grothendieck's algebraization theorem for
formal coherent sheaves (the existence theorem, \cite[III.5]{ega}),
$\fX'$ is the formal completion of an étale covering $\bar X'$ of
$\bar X=X\otimes_A\hat A$.

By passage to the limit, it suffices to prove the proposition when $A$ is
the henselization of a finitely generated $\dZ$-algebra. We can then
apply Artin's approximation theorem to the functor
$\sF:(\textnormal{$A$-algebras})\to(\textnormal{sets})$ that associates
to an $A$-algebra $B$ the set of isomorphism classes of étale coverings
of $X\otimes_AB$. Indeed, this functor is locally of finite presentation:
if $B_i$ is a filtered direct system of $A$-algebras and
$B=\varinjlim B_i$, then $\sF(B)=\varinjlim\sF(B_i)$. By Artin's
theorem, given an element $\bar\xi\in\sF(\hat A)$, here the isomorphism class
of $\bar X'$, there exists $\xi\in\sF(A)$ having the same image as
$\bar\xi$ in $\sF(A/\fm)$. In other words, there exists an étale
covering $X'$ of $X$ whose restriction to $X_0$ is isomorphic to
$X_0'$.

\subsection{Constructible sheaves}\label{I:4-3}

In this paragraph, $X$ is a \emph{noetherian} scheme, and by a sheaf on
$X$ we mean an \emph{abelian} sheaf on $\et X$.

\begin{definition}\label{I:4-3-1}
A sheaf $\sF$ on $X$ is called \emph{locally constant constructible}
(abbreviated l.c.c.) if it is represented by an étale covering of $X$.
\end{definition}

\begin{definition}\label{I:4-3-2}
A sheaf $\sF$ on $X$ is called \emph{constructible} if it satisfies the
following equivalent conditions:
\begin{enumerate}[\indent (i)]
  \item There exists a finite surjective family of subschemes $X_i$ of
    $X$ such that the restriction of $\sF$ to $X_i$ is l.c.c.
  \item There exist a finite family of finite morphisms $p_i:X_i'\to X$,
    for each $i$ a constant constructible sheaf $C_i$ (i.e. one defined
    by a finite abelian group) on $X_i'$, and a monomorphism
    $\sF\to\prod p_{i*}C_i$.
\end{enumerate}
\end{definition}

It is readily checked that the category of constructible sheaves on $X$
is an abelian category. Moreover, if $u:\sF\to\sG$ is a homomorphism of
sheaves and $\sF$ is constructible, then $\im(u)$ is constructible.

\begin{lemma}\label{I:4-3-3}
Every torsion sheaf $\sF$ is a filtered direct limit of constructible
sheaves.
\end{lemma}

Indeed, if $j:U\to X$ is an étale scheme of finite type over $X$, an
element $\xi\in\sF(U)$ such that $n\xi=0$ defines a homomorphism of
sheaves $j_!\const{\dZ/n}\to\sF$ whose image (the smallest subsheaf of
$\sF$ for which $\xi$ is a local section) is a constructible subsheaf
of $\sF$. Clearly $\sF$ is the direct limit of such subsheaves.

\begin{definition}\label{I:4-3-4}
Let $\cC$ be an abelian category and let $T$ be a functor from $\cC$ to
the category of abelian groups. The functor $T$ is called
\emph{effaceable} in $\cC$ if, for every object $A$ of $\cC$ and every
$\alpha\in T(A)$, there exists a monomorphism $u:A\to M$ in $\cC$ such
that $T(u)\alpha=0$.
\end{definition}

\begin{lemma}\label{I:4-3-5}
The functors $\h^q(X,-)$ for $q>0$ are effaceable in the category of
constructible sheaves on $X$.
\end{lemma}

It suffices to observe that, if $\sF$ is a constructible sheaf, there is
necessarily an integer $n>0$ such that $\sF$ is a sheaf of
$\dZ/n$-modules. There is then a monomorphism $\sF\hookrightarrow\sG$,
where $\sG$ is a sheaf of $\dZ/n$-modules and
$\h^q(X,\sG)=0$ for every $q>0$. For example, one can take for $\sG$
the Godement resolution $\prod_{x\in X}i_{x*}\sF_{\bar x}$, where $x$
runs over the points of $X$ and $i_x:\bar x\to X$ is a geometric point
centred at $x$. By (\ref{I:4-3-3}), $\sG$ is a direct limit of
constructible sheaves, whence the lemma, because the functors
$\h^q(X,-)$ commute with direct limits.

\begin{lemma}\label{I:4-3-6}
Let $\varphi^\bullet:T^\bullet\to{T'}^\bullet$ be a morphism of
cohomological functors defined on an abelian category $\cC$ with values
in the category of abelian groups. Suppose that $T^q$ is effaceable for
$q>0$, and let $\cE$ be a subset of the objects of $\cC$ such that
every object of $\cC$ is contained in an object belonging to $\cE$.
Then the following conditions are equivalent:
\begin{enumerate}[\indent (i)]
  \item $\varphi^q(A)$ is bijective for every $q\geqslant0$ and every
    $A\in\ob\cC$.
  \item $\varphi^0(M)$ is bijective and $\varphi^q(M)$ is surjective for
    every $q>0$ and every $M\in\cE$.
  \item $\varphi^0(A)$ is bijective for every $A\in\ob\cC$, and
    ${T'}^q$ is effaceable for every $q>0$.
\end{enumerate}
\end{lemma}

The proof is by induction on $q$ and presents no difficulty.

\begin{proposition}\label{I:4-3-7}
Let $X_0$ be a subscheme of $X$. Suppose that, for every
$n>0$ and every scheme $X'$ finite over $X$, the canonical map
\[
  \h^q(X',\dZ/n) \to \h^q(X_0',\dZ/n)\text{,}
\]
where $X_0'=X'\times_XX_0$, is bijective for $q=0$ and surjective for
$q>0$. Then, for every torsion sheaf $\sF$ on $X$ and every
$q\geqslant0$, the canonical map
\[
  \h^q(X,\sF) \to \h^q(X_0,\sF)
\]
is bijective.
\end{proposition}

By passage to the limit, it suffices to prove the assertion for
constructible $\sF$. Apply Lemma (\ref{I:4-3-6}) with $\cC$ the category
of constructible sheaves on $X$, $T^q=\h^q(X,-)$,
${T'}^q=\h^q(X_0,-)$, and $\cE$ the set of constructible sheaves of the
form $\prod p_{i*}C_i$, where $p_i:X_i'\to X$ is finite and $C_i$ is a
finite constant sheaf on $X_i'$.

\subsection{End of the proof}\label{I:4-4}

By the method of fibration by curves (\ref{I:3-4-3}), we reduce to proving
the theorem in relative dimension $\leqslant1$. By the preceding
paragraph, it suffices to show the following. If $S$ is the spectrum of a
strictly henselian noetherian local ring, $f:X\to S$ is a proper
morphism whose closed fibre $X_0$ has dimension $\leqslant1$, and $n$
is an integer $>0$, then the canonical homomorphism
\[
  \h^q(X,\dZ/n) \to \h^q(X_0,\dZ/n)
\]
is bijective for $q=0$ and surjective for $q>0$.

The cases $q=0$ and $1$ were treated above, and
$\h^q(X_0,\dZ/n)=0$ for $q\geqslant3$; it therefore suffices to treat
$q=2$. We may clearly suppose that $n$ is a power of a prime. If
$n=p^r$, where $p$ is the characteristic of the residue field of $S$,
Artin--Schreier theory shows that $\h^2(X_0,\dZ/p^r)=0$. If
$n=\ell^r$, $\ell\ne p$, Kummer theory gives a commutative diagram
\[\xymatrix{
  \pic(X) \ar[d] \ar[r]^-\alpha
    & \h^2(X,\dZ/\ell^r) \ar[d] \\
  \pic(X_0) \ar[r]^-\beta
    & \h^2(X_0,\dZ/\ell^r)
}\]
in which $\beta$ is surjective (this was seen in Chapter \ref{I:3} for a
smooth curve over an algebraically closed field, but similar arguments
apply to any curve over a separably closed field).

Thus it suffices to prove:

\begin{proposition}\label{I:4-4-1}
Let $S$ be the spectrum of a henselian noetherian local ring, and let
$f:X\to S$ be a proper morphism whose closed fibre $X_0$ has dimension
$\leqslant1$. Then the canonical restriction map
\[
  \pic(X)\to \pic(X_0)
\]
is surjective. (In fact, it is enough that $f$ be a \emph{separated}
morphism of finite type.)
\end{proposition}

To simplify the proof, we suppose that $X$ is integral, although this is
not necessary. Every invertible sheaf on $X_0$ is associated with a
Cartier divisor (because $X_0$ is a curve, hence quasi-projective), so it
suffices to show that the canonical map $\Div(X)\to\Div(X_0)$ is
surjective.

Every divisor on $X_0$ is a linear combination of divisors whose support
is concentrated at a single nonisolated closed point of $X_0$. Let $x$
be such a point, let $t_0\in\sO_{X_0,x}$ be a regular nonunit in
$\sO_{X_0,x}$, and let $D_0$ be the divisor concentrated at $x$ with
local equation $t_0$. Let $U$ be an open neighbourhood of $x$ in $X$
on which there is a section $t\in\Gamma(U,\sO_U)$ lifting $t_0$. Let
$Y$ be the closed subset of $U$ defined by $t=0$; after making $U$
sufficiently small, we may suppose that $x$ is the only point of
$Y\cap X_0$. Then $Y$ is quasi-finite over $S$ at $x$. Since $S$ is
the spectrum of a henselian local ring, it follows that
$Y=Y_1\amalg Y_2$, where $Y_1$ is finite over $S$ and $Y_2$ does not
meet $X_0$. Moreover, since $X$ is separated over $S$, $Y_1$ is closed
in $X$.

After replacing $U$ by a smaller open neighbourhood of $x$, we may
suppose that $Y=Y_1$, in other words, that $Y$ is closed in $X$. We then
define a divisor $D$ on $X$ lifting $D_0$ by setting
$D|X\setminus Y=0$ and $D|U=\dv(t)$; this makes sense because $t$ is
invertible on $U\setminus Y$.

\subsubsection{Remark}\label{I:4-4-2}

When $f$ is proper, one could also give a proof similar to that of
Proposition (\ref{I:4-2-2}). Indeed, since $X_0$ is a curve, there is no
obstruction to lifting an invertible sheaf on $X_0$ to the infinitesimal
neighbourhoods $X_n$ of $X_0$, and hence to the formal completion $\fX$
of $X$ along $X_0$. One then concludes by applying successively
Grothendieck's existence theorem and Artin's approximation theorem.

\subsection{Cohomology with proper support}\label{I:4-5}

\begin{definition}\label{I:4-5-1}
Let $X$ be a separated scheme of finite type over a field $k$. By a
theorem of Nagata, there exist a scheme $\bar X$ proper over $k$ and an
open immersion $j:X\to\bar X$. For every torsion sheaf $\sF$ on $X$,
let $j_!\sF$ denote the extension by $0$ of $\sF$ to $\bar X$, and
define the groups of \emph{cohomology with proper support} by
\[
  \h_c^q(X,\sF)=\h^q(\bar X, j_! \sF) \text{.}
\]
\end{definition}

Let us show that this definition is independent of the chosen
compactification $j:X\to\bar X$. Let $j_1:X\to\bar X_1$ and
$j_2:X\to\bar X_2$ be two compactifications. Then $X$ maps into
$\bar X_1\times\bar X_2$ by $x\mapsto(j_1(x),j_2(x))$, and
the closed image $\bar X_3$ of $X$ under this map is a compactification
of $X$. We thus have a commutative diagram
\[\xymatrix{
  & \bar X_1 \\
  X \ar[ur]^-{j_1} \ar[rr]^-{(j_1,j_2)} \ar[dr]_-{j_2}
    & & \bar X_3 \ar[ul]_-{p_1} \ar[dl]^-{p_2} \\
  & \bar X_2
}\]
where $p_1$ and $p_2$, the restrictions to $\bar X_3$ of the natural
projections, are proper morphisms. It therefore suffices to treat the case
of a commutative diagram
\[\xymatrix{
  X \ar[r]^-{j_2} \ar[dr]_-{j_1}
    & \bar X_2 \ar[d]^-p \\
  & \bar X_1
}\]
with $p$ proper.

\begin{lemma}\label{I:4-5-2}
One has $p_*(j_{2!}\sF)=j_{1!}\sF$ and
$\eR^qp_*(j_{2!}\sF)=0$ for $q>0$.
\end{lemma}

Let us note at once that the lemma suffices. Using the Leray spectral
sequence for $p$, it gives, for every $q\geqslant0$,
\[
  \h^q(\bar X_2, j_{2!} \sF) = \h^q(\bar X_1, j_{1!} \sF) \text{.}
\]

To prove the lemma, argue fibre by fibre using the base change theorem
(\ref{I:4-1-1}) for $p$. The result is immediate because above a point
of $X$, $p$ is an isomorphism, whereas above a point of
$\bar X_1\setminus X$, $j_{2!}\sF$ vanishes on the fibre of $p$.

\subsubsection{}\label{I:4-5-3}

Similarly, if $f:X\to S$ is a separated morphism of finite type between
noetherian schemes, there exist a proper morphism $\bar f:\bar X\to S$
and an open immersion $j:X\to\bar X$. The higher direct images with
proper support $\eR^qf_!$ are defined, for every torsion sheaf $\sF$ on
$X$, by
\[
  \eR^q f_! \sF = \eR^q \bar f_*(j_! \sF) \text{.}
\]
As before, one checks that this definition is independent of the chosen
compactification.

\begin{theorem}\label{I:4-5-4}
Let $f:X\to S$ be a separated morphism of finite type between noetherian
schemes, and let $\sF$ be a torsion sheaf on $X$. Then the fibre of
$\eR^qf_!\sF$ at a geometric point $s$ of $S$ is isomorphic to the
cohomology with proper support $\h_c^q(X_s,\sF)$ of the fibre $X_s$ of
$f$ at $s$.
\end{theorem}

This is a simple variant of the proper base change theorem
(\ref{I:4-1-1}). More generally, if
\[\xymatrix{
  X \ar[d]^-f
    & X' \ar[l]_-{g'} \ar[d]^-{f'} \\
  S
    & S' \ar[l]^-g
}\]
is a cartesian diagram, there is a canonical isomorphism
\begin{equation*}\tag{4.5.4.1}\label{I:eq:4-5-4-1}
  g^* (\eR^q f_! \sF) \simeq \eR^q f'_!({g'}^* \sF) \text{.}
\end{equation*}

\subsection{Applications}\label{I:4-6}

\begin{theorem}[Vanishing theorem]\label{I:4-6-1}
Let $f:X\to S$ be a separated morphism of finite type whose fibres have
dimension $\leqslant n$, and let $\sF$ be a torsion sheaf on $X$. Then
$\eR^qf_!\sF=0$ for $q>2n$.
\end{theorem}

By the base change theorem, we may suppose that $S$ is the spectrum of a
separably closed field. If $\dim X=n$, there exists an affine open
$U$ of $X$ such that $\dim(X\setminus U)<n$. We then have an exact
sequence $0\to\sF_U\to\sF\to\sF_{X\setminus U}\to0$, and induction on
$n$ reduces us to proving the theorem for $X=U$ affine. The method of
fibration by curves (\ref{I:3-4-1}) and the base change theorem then
reduce us to a curve over a separably closed field, for which the desired
result follows from Tsen's theorem (\ref{I:3-3-6}).

\begin{theorem}[Finiteness theorem]\label{I:4-6-2}
Let $f:X\to S$ be a separated morphism of finite type and let $\sF$ be
a constructible sheaf on $X$. Then the sheaves $\eR^qf_!\sF$ are
constructible.
\end{theorem}

We consider only the case in which $\sF$ is annihilated by an integer
invertible on $X$.

To prove the theorem, one reduces to the case in which $\sF$ is the
constant sheaf $\const{\dZ/n}$ and $f:X\to S$ is proper and smooth with
geometrically connected curve fibres of genus $g$. If $n$ is invertible
on $X$, the sheaves $\eR^qf_*\sF$ are then locally free of finite rank
and vanish for $q>2$ (\ref{I:4-6-1}). Replacing $\dZ/n$ by the sheaf
$\dmu_n$, which is locally isomorphic to it (over $S$), there are
canonical isomorphisms
\begin{equation*}\tag{4.6.2.1}\label{I:eq:4-6-2-1}
\begin{aligned}
  \eR^0 f_* \dmu_n &= \dmu_n \\
  \eR^1 f_* \dmu_n &= \const\pic(X/S)_n \\
  \eR^2 f_* \dmu_n &= \dZ/n \text{.}
\end{aligned}
\end{equation*}

\begin{theorem}[Comparison with classical cohomology]\label{I:4-6-3}
Let $f:X\to S$ be a separated morphism of schemes of finite type over
$\dC$, and let $\sF$ be a torsion sheaf on $X$. Denote by a superscript
$\an{(-)}$ the functor passing to the usual topological spaces, and by
$\eR^q\an f_!$ the derived functors of direct image with proper support
under $\an f$. Then
\[
  \an{(\eR^q f_! \sF)} \simeq \eR^q \an f_! \an \sF \text{.}
\]
In particular, when $S=$ a point and $\sF$ is the constant sheaf
$\dZ/n$,
\[
  \h_c^q(X,\dZ/n) \simeq \h_c^q(\an X,\dZ/n)\text{.}
\]
\end{theorem}

Unwindings using the base change theorem reduce us to the case in which
$X$ is a smooth proper curve, $S=$ a point, and $\sF=\dZ/n$. The
cohomology groups in question then vanish for $q\ne0,1,2$, and we invoke
\cite{se55}: indeed, if $X$ is proper over $\dC$, then
$\pi_0(X)=\pi_0(\an X)$ and $\pi_1(X)=$ the profinite completion of
$\pi_1(\an X)$, which proves the assertion for $q=0,1$. For $q=2$,
one uses the Kummer exact sequence and the fact that, again by
\cite{se55}, $\pic(X)=\pic(\an X)$.

\begin{theorem}[Cohomological dimension of affine schemes]\label{I:4-6-4}
Let $X$ be an affine scheme of finite type over a separably closed field
and let $\sF$ be a torsion sheaf on $X$. Then
$\h^q(X,\sF)=0$ for $q>\dim(X)$.
\end{theorem}

For the very elegant proof, we refer to \cite{sga4}, XIV \S 2 and 3.

\subsubsection{Remark}\label{I:4-6-5}

This theorem is in a sense a substitute for Morse theory. Indeed, consider
the classical case in which $X$ is smooth and affine over $\dC$, embedded
in a standard affine space $\dC^N$. Then, for almost every point
$p\in\dC^N$, the function “distance from $p$” on $X$ is a Morse
function, and the indices of its critical points are smaller than
$\dim(X)$. Thus $X$ is obtained by attaching handles of index less than
$\dim(X)$, which gives the classical analogue of (\ref{I:4-6-4}).

\section{Local acyclicity of smooth morphisms}\label{I:5}

Let $X$ be a complex analytic variety and let $f:X\to D$ be a morphism
from $X$ to the disk. Write $[0,t]$ for the closed line segment in $D$
with endpoints $0$ and $t$, and $(0,t]$ for the half-open segment. If
$f$ is \emph{smooth}, the inclusion
\[
  j : f^{-1}\left((0,t]\right) \hookrightarrow f^{-1}\left([0,t]\right)
\]
is a homotopy equivalence; the special fibre
$X_0=f^{-1}(0)$ can be pushed into $f^{-1}([0,t])$.

In practice, for sufficiently small $t$, $f^{-1}((0,t])$ will be a
bundle over $(0,t]$, so that the inclusion
\[
  X_t = f^{-1}(t) \hookrightarrow f^{-1}\left((0,t]\right)
\]
will also be a homotopy equivalence. The \emph{cospecialization morphism}
is then the homotopy class of maps
\[
  \cosp : X_0 \hookrightarrow f^{-1}\left([0,t]\right) \xleftarrow\sim f^{-1}\left((0,t]\right) \xleftarrow\sim X_t \text{.}
\]
Figuratively speaking,
for a smooth morphism the general fibre swallows the special fibre.

No longer suppose that $f$ is necessarily smooth (but suppose that
$f^{-1}((0,t])$ is a bundle over $(0,t]$). A morphism
$\cosp^\bullet$ can still be defined in cohomology whenever
$j_*\dZ=\dZ$ and $\eR^qj_*\dZ=0$ for $q>0$. Under these hypotheses,
the Leray spectral sequence for $j$ shows that
\[
  \h^\bullet\left(f^{-1}\left([0,t]\right),\dZ\right) \iso \h^\bullet\left(f^{-1}\left((0,t]\right),\dZ\right)
\]
and $\cosp^\bullet$ is the composite morphism
\[
  \cosp^\bullet : \h^\bullet(X_t,\dZ)
    \xleftarrow\sim \h^\bullet\left(f^{-1}\left((0,t]\right),\dZ\right)
    \xleftarrow\sim \h^\bullet\left(f^{-1}\left([0,t]\right),\dZ\right)
    \to \h^\bullet(X_0,\dZ)
\]
The fibre of $\eR^qj_*\dZ$ at a point $x\in X_0$ is calculated as
follows. In an ambient space, take a ball $B_\varepsilon$ centred at $x$
with sufficiently small radius $\varepsilon$, and for sufficiently small
$\eta$ put $E=X\cap B_\varepsilon\cap f^{-1}(\eta t)$; this is the
\emph{vanishing-cycle variety} at $x$. One has
\[
  (\eR^q j_* \dZ)_x \xleftarrow\sim \h^q\left(X\cap B_\varepsilon\cap f^{-1}\left((0,\eta t]\right),\dZ\right) \xrightarrow\sim \h^q(E,\dZ)
\]
and the cospecialization morphism is defined in cohomology whenever the
vanishing-cycle varieties are acyclic
($\h^0(E,\dZ)=\dZ$ and $\h^q(E,\dZ)=0$ for $q>0$), which is expressed
by saying that $f$ is \emph{locally acyclic}.

This chapter is devoted to the analogue of this situation for a smooth
morphism of schemes in étale cohomology. In this setting, however, it is
essential to restrict to torsion coefficients of order \emph{prime to the
residue characteristics}. Paragraph \ref{I:5-1} gives generalities on
locally acyclic morphisms and cospecialization maps. In Paragraph
\ref{I:5-2}, we prove that a smooth morphism is locally acyclic. In
Paragraph \ref{I:5-3}, this result is combined with those of the preceding
chapter to obtain two applications: a specialization theorem for
cohomology groups (the cohomology of the geometric fibres of a smooth
proper morphism is locally constant), and a base change theorem for a
smooth morphism.

Throughout what follows, we fix an integer $n$, and “scheme” means
“scheme on which $n$ is invertible.” A “geometric point” will always mean
an “algebraic geometric point” (\ref{I:2-3-1})
$x:\spec(k)\to X$, with $k$ algebraically closed.

\subsection{Locally acyclic morphisms}\label{I:5-1}

\subsubsection{Notation}\label{I:5-1-1}

Given a scheme $S$ and a geometric point $s$ of $S$, we denote by
$\widetilde S^s$ the spectrum of the strict localization of $S$ at $s$.

\begin{definition}\label{I:5-1-2}
A geometric point $t$ of $S$ is called a \emph{generization} of $s$ if
it is defined by an algebraic closure of the residue field of a point of
$\widetilde S^s$. We also say that $s$ is a \emph{specialization} of
$t$, and call the \emph{specialization map} the $S$-morphism
$t\to\widetilde S^s$.
\end{definition}

\begin{definition}\label{I:5-1-3}
Let $f:X\to S$ be a morphism of schemes. Let $s$ be a geometric point
of $S$, $t$ a generization of $s$, $x$ a geometric point of $X$ over
$s$, and
$\widetilde X_t^x=\widetilde X^x\times_{\widetilde S^s}t$. Then
$\widetilde X_t^x$ is called a \emph{vanishing-cycle variety} of $f$ at
$x$.

The morphism $f$ is called \emph{locally acyclic} if the reduced
cohomology of every vanishing-cycle variety $\widetilde X_t^x$ vanishes:
\begin{equation*}\tag{5.1.3.1}\label{I:eq:5-1-3-1}
  \hat\h^\bullet(\widetilde X_t^x,\dZ/n) = 0\text{,}
\end{equation*}
i.e. $\h^0\left(\widetilde X_t^x,\dZ/n\right)=\dZ/n$ and
$\h^q(\widetilde X_t^x,\dZ/n)=0$ for $q>0$.
\end{definition}

\begin{lemma}\label{I:5-1-4}
Let $f:X\to S$ be a locally acyclic morphism and let $g:S'\to S$ be a
quasi-finite morphism (or a projective limit of quasi-finite morphisms).
Then the morphism $f':X'\to S'$ obtained from $f$ by base change is
locally acyclic.
\end{lemma}

Indeed, every vanishing-cycle variety of $f'$ is a vanishing-cycle
variety of $f$.

\begin{lemma}\label{I:5-1-5}
Let $f:X\to S$ be a locally acyclic morphism. For every geometric point
$t$ of $S$ giving rise to a cartesian diagram
\[\xymatrix{
  X_t \ar[r]^-{\varepsilon'} \ar[d]
    & X \ar[d]^-f \\
  t \ar[r]^-\varepsilon
    & S
}\]
one has $\varepsilon_*'\dZ/n=f^*\varepsilon_*\dZ/n$ and
$\eR^q\varepsilon_*'\dZ/n=0$ for $q>0$.
\end{lemma}

Let $\bar S$ be the closure of $\varepsilon(t)$, let $S'$ be the
normalization of $\bar S$ in $k(t)$, and consider the cartesian diagram
\[\xymatrix{
  X_t \ar[r]^-{i'} \ar[d]
    & X' \ar[r]^-{\alpha'} \ar[d]^-{f'}
    & X \ar[d]^-f \\
  t \ar[r]^-i
    & S' \ar[r]^-\alpha
    & S
}\]
The local rings of $S'$ are normal with separably closed fraction fields.
They are therefore strictly henselian, and the local acyclicity of $f'$
(\ref{I:5-1-4}) gives
$i'_*\dZ/n={f'}^*i_*\dZ/n$ and $\eR^q i'_*\dZ/n=0$ for $q>0$. Since $\alpha$
is integral, it follows that
\[
  \eR^q\varepsilon_*'\dZ/n
    = \alpha_*' \eR^q i_*' \dZ/n
    = \alpha_*' {f'}^* \eR^q i_* \dZ/n
    = f^* \alpha_* \eR^q i_* \dZ/n
    = f^* \eR^q \varepsilon_* \dZ/n \text{,}
\]
which proves the lemma.

\subsubsection{}\label{I:5-1-6}

Given a locally acyclic morphism $f:X\to S$ and a specialization map
$t\to\widetilde S^s$, we shall define canonical homomorphisms, called
cospecialization maps,
\[
  \cosp^\bullet : \h^\bullet(X_t,\dZ/n) \to \h^\bullet(X_s,\dZ/n)\text{,}
\]
relating the cohomology of the general fibre $X_t=X\times_St$ to that of
the special fibre $X_s=X\times_Ss$.

Consider the cartesian diagram
\[\xymatrix{
  X_t \ar[d] \ar[r]^-{\varepsilon'}
    & \widetilde X \ar[d]^-{f'}
    & \ar[l] X_s \ar[d] \\
  t \ar[r]^-\varepsilon
    & \widetilde S^s
    & \ar[l] s
}\]
obtained from $f$ by base change. By (\ref{I:5-1-4}), $f'$ is again
locally acyclic. The definition of local acyclicity immediately shows that
the restrictions to $X_s$ of the sheaves $\eR^q\varepsilon_*'\dZ/n$ are
$\dZ/n$ for $q=0$ and $0$ for $q>0$. By (\ref{I:5-1-5}), we even know
that $\eR^q\varepsilon_*'\dZ/n=0$ for $q>0$. We define
$\cosp^\bullet$ as the composite map
\begin{equation*}\tag{5.1.6.1}\label{I:eq:5-1-6-1}
  \h^\bullet(X_t,\dZ/n) \simeq \h^\bullet(\widetilde X,\varepsilon_*' \dZ/n) \to \h^\bullet(X_s,\dZ/n) \text{.}
\end{equation*}

\paragraph{Variant}
Let $\bar S$ be the closure of $\varepsilon(t)$ in
$\widetilde S^s$, let $S'$ be the normalization of $\bar S$ in $k(t)$,
and let $X'/S'$ be obtained from $X/S$ by base change. Diagram
\eqref{I:eq:5-1-6-1} may also be written
\[
  \h^\bullet(X_t,\dZ/n) \simeq \h^\bullet(X',\dZ/n) \to \h^\bullet(X_s,\dZ/n) \text{.}
\]

\begin{theorem}\label{I:5-1-7}
Let $S$ be a locally noetherian scheme, $s$ a geometric point of $S$,
and $f:X\to S$ a morphism. Suppose that
\begin{enumerate}[\indent a)]
  \item $f$ is locally acyclic,
  \item for every specialization map $t\to\widetilde S^s$ and every
    $q\geqslant0$, the cospecialization maps
    $\h^q(X_t,\dZ/n)\to\h^q(X_s,\dZ/n)$ are bijective.
\end{enumerate}
Then the canonical homomorphism
$(\eR^qf_*\dZ/n)_s\to\h^q(X_s,\dZ/n)$ is bijective for every
$q\geqslant0$.
\end{theorem}

To prove the theorem, we may clearly suppose that
$S=\widetilde S^s$. We shall in fact show that, for every sheaf of
$\dZ/n\dZ$-modules $\sF$ on $S$, the canonical homomorphism
$\varphi^q(\sF):(\eR^qf_*f^*\sF)_s\to\h^q(X_s,f^*\sF)$ is bijective.

Every sheaf of $\dZ/n\dZ$-modules is a filtered direct limit of
constructible sheaves of $\dZ/n\dZ$-modules (\ref{I:4-3-3}). Moreover,
every constructible sheaf of $\dZ/n\dZ$-modules embeds in a sheaf of the
form $\prod i_{\lambda *}C_\lambda$, where
$i_\lambda:t_\lambda\to S$ is a finite family of generizations of $s$
and $C_\lambda$ is a free $\dZ/n\dZ$-module of finite rank on
$t_\lambda$. By the definition of the cospecialization maps, condition
b) says that the homomorphisms $\varphi^q(\sF)$ are bijective when
$\sF$ has this form.

We conclude by a variant of Lemma (\ref{I:4-3-6}):

\begin{lemma}\label{I:5-1-8}
Let $\cC$ be an abelian category in which filtered direct limits exist.
Let $\varphi^\bullet:T^\bullet\to{T'}^\bullet$ be a morphism of
cohomological functors commuting with filtered direct limits, defined on
$\cC$ with values in the category of abelian groups. Suppose that there
are two subsets $\cD$ and $\cE$ of the objects of $\cC$ such that:
\begin{enumerate}[\indent a)]
  \item every object of $\cC$ is a filtered direct limit of objects
    belonging to $\cD$,
  \item every object belonging to $\cD$ is contained in an object
    belonging to $\cE$.
\end{enumerate}
Then the following conditions are equivalent:
\begin{enumerate}[\indent (i)]
  \item $\varphi^q(A)$ is bijective for every $q\geqslant0$ and every
    $A\in\ob\cC$.
  \item $\varphi^q(M)$ is bijective for every $q\geqslant0$ and every
    $M\in\cE$.
\end{enumerate}
\end{lemma}

The proof of the lemma proceeds by passage to the direct limit, induction
on $q$, and repeated application of the five lemma to the diagram of
exact cohomology sequences deduced from an exact sequence
$0\to A\to M\to A'\to0$, with $A\in\cD$, $M\in\cE$, and
$A'\in\ob\cC$.

\begin{corollary}\label{I:5-1-9}
Let $S$ be the spectrum of a strictly henselian noetherian local ring and
let $f:X\to S$ be a locally acyclic morphism. Suppose that for every
geometric point $t$ of $S$ one has $\h^0(X_t,\dZ/n)=\dZ/n$ and
$\h^q(X_t,\dZ/n)=0$ for $q>0$ (in other words, the geometric fibres of
$f$ are acyclic). Then $f_*\dZ/n=\dZ/n$ and
$\eR^qf_*\dZ/n=0$ for $q>0$.
\end{corollary}

\begin{corollary}\label{I:5-1-10}
Let $f:X\to Y$ and $g:Y\to Z$ be morphisms of locally noetherian
schemes. If $f$ and $g$ are locally acyclic, then so is $g\circ f$.
\end{corollary}

We may suppose that $X$, $Y$, and $Z$ are strictly local and that $f$
and $g$ are local morphisms. It is then enough to show that, if $z$ is
an algebraic geometric point of $Z$, then
$\h^0(X_z,\dZ/n)=\dZ/n$ and $\h^q(X_z,\dZ/n)=0$ for $q>0$.

Since $g$ is locally acyclic, one has $\h^0(Y_z,\dZ/n)=\dZ/n$ and
$\h^q(Y_z,\dZ/n)=0$ for $q>0$. Moreover, the morphism
$f_z:X_z\to Y_z$ is locally acyclic (\ref{I:5-1-4}), and its geometric
fibres are acyclic because they are vanishing-cycle varieties of $f$.
By (\ref{I:5-1-9}), $\eR^q{f_z}_*\dZ/n=0$ for $q>0$. In addition,
${f_z}_*\dZ/n$ is constant with fibre $\dZ/n$ on $Y_z$. The result
follows from the Leray spectral sequence for $f_z$.

\subsection{Local acyclicity of a smooth morphism}\label{I:5-2}

\begin{theorem}\label{I:5-2-1}
A smooth morphism is locally acyclic.
\end{theorem}

Let $f:X\to S$ be a smooth morphism. The assertion is local for the étale
topology on $X$ and $S$, so we may suppose that $X$ is standard affine
space of dimension $d$ over $S$. By passage to the limit, we may suppose
that $S$ is noetherian, and the transitivity of local acyclicity
(\ref{I:5-1-10}) shows that it suffices to treat $d=1$.

Let $s$ be a geometric point of $S$ and $x$ a geometric point of $X$
centred at a closed point of $X_s$. We must show that the geometric fibres
of $\widetilde X^x\to\widetilde S^s$ are acyclic. Henceforth put
$S=\widetilde S^s=\spec(A)$ and $X=\widetilde X^x$. Then
$X\simeq\spec A\{T\}$, where $A\{T\}$ is the henselization of $A[T]$
at the point $T=0$ above $s$.

If $t$ is a geometric point of $S$, the fibre $X_t$ is a projective
limit of smooth affine curves over $t$. Thus $\h^q(X_t,\dZ/n)=0$ for
$q\geqslant2$, and it suffices to show that
$\h^0(X_t,\dZ/n)=\dZ/n$ and $\h^1(X_t,\dZ/n)=0$ for $n$ prime to the
residue characteristic of $S$. This follows from the next two
propositions.

\begin{proposition}\label{I:5-2-2}
Let $A$ be a strictly henselian local ring, $S=\spec(A)$, and
$X=\spec A\{T\}$. Then the geometric fibres of $X\to S$ are connected.
\end{proposition}

By passage to the limit, we may reduce to the case in which $A$ is a
strict henselization of a finitely generated $\dZ$-algebra.

Let $\bar t$ be a geometric point of $S$ lying over $t$, and let $k'$
be a finite separable extension of $k(t)$ in $k(\bar t)$. Put
$t'=\spec(k')$ and $X_{t'}=X\times_S\spec(k')$. We must check that, for
all $\bar t$ and $t'$, $X_{t'}$ is connected (by which we mean connected
and nonempty). Let $A'$ be the normalization of $A$ in $k'$, i.e. the
ring of elements of $k'$ integral over the image of $A$ in $k(t)$. One
has $A\{T\}\otimes_AA'\iso A'\{T\}$: indeed, the left-hand side is
henselian local (because $A'$ is finite over $A$ and local) and is a
limit of local étale algebras over
$A'\{T\}=A\{T\}\otimes_AA'$. Thus $X_{t'}$ is still the fibre at $t'$
of $X=\spec\left(A'\{T\}\right)$ over $S'=\spec(A')$. The local scheme
$X'$ is normal, hence integral; its localization $X_{t'}$ is still
integral, and a fortiori connected.

\begin{proposition}\label{I:5-2-3}
Let $A$ be a strictly henselian local ring, $S=\spec(A)$, and
$X=\spec\left(A\{T\}\right)$. Let $\bar t$ be a geometric point of $S$
and $X_{\bar t}$ the corresponding geometric fibre. Then every Galois
étale covering of $X_{\bar t}$ of order prime to the characteristic of
the residue field of $A$ is trivial.
\end{proposition}

\begin{lemma}[Zariski--Nagata purity theorem in dimension $2$] \label{I:5-2-4} % originally 5.2.3.1
Let $C$ be a regular local ring of dimension $2$ and let $C'$ be a
finite normal $C$-algebra étale over the open complement of the closed
point of $\spec(C)$. Then $C'$ is étale over $C$.
\end{lemma}

Indeed, $C'$ is normal of dimension $2$, and hence
$\operatorname{prof}(C')=2$. Since
$\operatorname{prof}(C')+\dim\operatorname{proj}(C')=\dim(C)=2$, it
follows that $C'$ is free over $C$. The set of points of $C$ where $C'$
is ramified is then defined by an equation, the discriminant; since it
contains no point of height $1$, it is empty.

\begin{lemma}[Special case of Abhyankar's lemma]\label{I:5-2-5} % originally 5.2.3.2
Let $S=\spec(V)$ be a trait, let $\pi$ be a uniformizer, let $\eta$ be
the generic point of $S$, let $X$ be smooth over $S$, irreducible, and
of relative dimension $1$, let $\widetilde X_\eta$ be a Galois étale
covering of $X_\eta$ of degree $n$ invertible on $S$, and let
$S_1=\spec\left(V[\pi^{1/n}]\right)$. Use a subscript $(-)_1$ to denote
base change from $S$ to $S_1$. Then $\widetilde X_{1\eta}$ extends to an
étale covering of $X_1$.
\end{lemma}

Let $\widetilde X_1$ be the normalization of $X_1$ in
$\widetilde X_{1\eta}$. By the structure of the tame inertia groups of
the discrete valuation rings obtained by localizing $X$ at the generic
points of the special fibre $X_s$, $\widetilde X_1$ is étale over $X_1$
on the general fibre and at the generic points of the special fibre. By
(\ref{I:5-2-4}), it is étale everywhere.

\subsubsection{}\label{I:5-2-6} % originally 5.2.3.3

Let $t$ denote the point over which $\bar t$ lies. We are free to replace
$A$ by the normalization of $A$ in a finite separable extension $k(t')$ of
$k(t)$ in $k(\bar t)$ (cf. \ref{I:5-2-2}). This and a preliminary
passage to the limit allow us to suppose that
\begin{enumerate}[\indent a)]
  \item $A$ is normal and noetherian, and $t$ is the generic point of $S$.
  \item The étale covering under consideration of $X_{\bar t}$ comes from
    an étale covering of $X_t$.
  \item It comes from an étale covering $\widetilde X_U$ of the inverse
    image $X_U$ of a nonempty open subset $U$ of $S$ (this follows from
    b): $t$ is the limit of the $U$).
  \item The complement of $U$ has codimension $\geqslant2$ (at the price
    of enlarging $k(t)$, by applying (\ref{I:5-2-5}) to the discrete
    valuation rings obtained by localizing $S$ at the points of
    $S\setminus U$ of codimension $1$ in $S$; there are only finitely many
    such points).
  \item The covering $\widetilde X_U$ is trivial over the subscheme $T=0$.
    (This may require enlarging $k(t)$ once more.)
\end{enumerate}

Under these conditions, we shall show that the covering
$\widetilde X_U$ is trivial.

\begin{lemma}\label{I:5-2-7} % originally 5.2.3.4
Let $A$ be a normal noetherian strictly henselian local ring, let $U$
be an open subset of $\spec(A)$ whose complement has codimension
$\geqslant2$, let $V$ be its inverse image in
$X=\spec\left(A\{T\}\right)$, and let $V'$ be an étale covering of
$V$. If $V'$ is trivial over $T=0$, then $V'$ is trivial.
\end{lemma}

Let $B=\Gamma(V',\sO)$. Since $V'$ is the inverse image of $V$ in
$\spec(B)$, it suffices to show that $B$ is finite étale over $A\{T\}$
(and hence split, since $A\{T\}$ is strictly henselian). Let
$\widehat X=\spec A\llbracket T\rrbracket$, and denote by $(-)^\wedge$ base
change from $X$ to $\widehat X$. The scheme $\widehat X$ is faithfully
flat over $X$. Hence
$\Gamma(\widehat V',\sO)=B\otimes_{A\{T\}}A\llbracket T\rrbracket$,
and it suffices to show that this ring $\widehat B$ is finite étale over
$A\llbracket T\rrbracket$.

Let $V_m$ (resp. $V_m'$) be the subscheme of $\widehat V$ (resp.
$\widehat V'$) defined by $T^{m+1}=0$. By hypothesis, $V_0'$ is a
trivial covering of $V_0$: a sum of $n$ copies of $V_0$. The same is true
of $V_m'/V_m$, because étale coverings are insensitive to nilpotents.
We obtain
\[
  \varphi : \Gamma(\widehat V',\sO) \to \varprojlim_m \Gamma(V_m',\sO) = \left(\varprojlim_m \Gamma(V_m,\sO)\right)^n \text{.}
\]
By hypothesis, the complement of $U$ has depth $\geqslant2$: one has
$\Gamma(V_m,\sO)=A[T]/(T^{m+1})$, and $\varphi$ is a homomorphism from
$\widehat B$ to $A\llbracket T\rrbracket^n$. Over $U$, it provides $n$
distinct sections of $\widehat V'/\widehat V$; hence $\widehat V'$ is
trivial, a sum of $n$ copies of $\widehat V$. Since the complement of
$\widehat V$ in $\widehat X$ still has codimension $\geqslant2$ (and
therefore depth $\geqslant2$), it follows that
$\widehat B=A\llbracket T\rrbracket^n$, proving the lemma.

\subsection{Applications}\label{I:5-3}

\begin{theorem}[Specialization of cohomology groups]\label{I:5-3-1}
Let $f:X\to S$ be a proper locally acyclic morphism, for example a smooth
proper morphism. Then the sheaves $\eR^qf_*\dZ/n$ are locally constant
constructible, and for every specialization map $t\to\widetilde S^s$,
the cospecialization maps
$\h^q(X_t,\dZ/n)\to\h^q(X_s,\dZ/n)$ are bijective.
\end{theorem}

This follows immediately from the definition of the cospecialization maps
and the finiteness and proper base change theorems.

\begin{theorem}[Base change by a smooth morphism]\label{I:5-3-2}
Let
\[\xymatrix{
  X' \ar[r]^-{g'} \ar[d]^-{f'}
    & X \ar[d]^-f \\
  S' \ar[r]^-{g}
    & S
}\]
be a cartesian diagram with $g$ smooth. For every sheaf $\sF$ on $X$
that is torsion of order prime to the residue characteristics of $S$,
one has
\[
  g^* \eR^q f_* \sF \iso \eR^q f'_*({g'}^* \sF) \text{.}
\]
\end{theorem}

After passing to an open covering of $X$, we reduce to the case where
$X$ is affine, and then, by passage to the limit, to the case where $X$
is of finite type over $S$. The morphism $f$ then factors as an open
immersion $j:X\to\bar X$ followed by a proper morphism
$\bar f:\bar X\to S$. The Leray spectral sequence for $\bar f\circ j$
and the proper base change theorem show that it suffices to prove the
theorem when $X\to S$ is an open immersion.

In this case, if $\sF$ has the form $\varepsilon_*C$, where
$\varepsilon:t\to X$ is a geometric point of $X$, the theorem is a
corollary of (\ref{I:5-1-5}). The general case follows from Lemma
(\ref{I:5-1-8}).

\begin{corollary}\label{I:5-3-3}
Let $K/k$ be an extension of separably closed fields, let $X$ be a
$k$-scheme, and let $n$ be an integer prime to the characteristic of
$k$. Then the canonical map
$\h^q(X,\dZ/n)\to\h^q(X_K,\dZ/n)$ is bijective for every
$q\geqslant0$.
\end{corollary}

It suffices to observe that $\bar K$ is a direct limit of smooth
$\bar k$-algebras.

\begin{theorem}[Relative purity]\label{I:5-3-4}
Let
\begin{equation*}\tag{5.3.4.1}\label{I:eq:5-3-4-1}\xymatrix{
  U \ar@{^{(}->}[r]^-j \ar[dr]
    & X \ar[d]^-f
    & \ar@{_{(}->}[l]_-i Y \ar[dl]^-h \\
  & S
}\end{equation*}
be a commutative diagram with $f$ smooth of pure relative dimension $N$,
$h$ smooth of pure relative dimension $N-1$, $i$ a closed immersion,
and $U=X\setminus Y$. If $n$ is prime to the residue characteristics of
$S$, then
\begin{equation*}
\begin{aligned}
  j_* \dZ/n     &= \dZ/n \\
  \eR^1 j_*\dZ/n &= \dZ/n(-1)_Y \\
  \eR^q j_*\dZ/n &= 0 \qquad \text{for $q\geqslant 2$}
\end{aligned}
\end{equation*}
\end{theorem}

In these formulas, $\dZ/n(-1)$ denotes the $\dZ/n$-dual of $\dmu_n$.
If $t$ is a local equation for $Y$, the isomorphism
$\eR^1j_*\dZ/n\simeq\dZ/n(-1)_Y$ is defined by the map
$a:\dZ/n\to\eR^1j_*\dmu_n$ sending $1$ to the class of the
$\dmu_n$-torsor of $n$th roots of $t$.

The question is local. This allows us to replace $(X,Y)$ by a locally
isomorphic pair, for example
\[\xymatrix{
  \dA_T^1 \ar@{^{(}->}[r]^-j \ar[dr]_-g
    & \dP_T^1 \ar[d]^-f
    & \ar@{_{(}->}[l]_-i T \ar[dl] \\
  & T
}\]
with $T=\dA_S^{N-1}$ and $i=$ the section at infinity. Corollary
(\ref{I:5-1-9}) applies to $g$ and gives
$\eR^qg_*\dZ/n=\dZ/n$ for $q=0$ and $0$ for $q>0$. For $f$, on the
other hand, \eqref{I:eq:4-6-2-1} gives
\[
  \eR^q f_* \dZ/n =
  \begin{cases}
    \dZ/n & q=0,\\
    0 & q=1,\\
    \dZ/n(-1) & q=2,\\
    0 & q>2.
  \end{cases}
\]
It is readily checked that $j_*\dZ/n=\dZ/n$, and that the
$\eR^qj_*\dZ/n=0$ are concentrated on $i(T)$ for $q>0$. The Leray
spectral sequence
$E_2^{pq}=\eR^pf_*\eR^qj_*\dZ/n\Rightarrow\eR^{p+q}g_*\dZ/n$ therefore
reduces to
\[\xymatrix{
  i^* \eR^q j_* \dZ/n & 0 & \cdots \\
  i^* \eR^1 j_* \dZ/n \ar[drr]^-{d_2} & 0 & \cdots \\
  \dZ/n              & 0 & \dZ/n(-1) & 0 & \cdots
}\]
$\eR^qj_*\dZ/n=0$ for $q\geqslant2$, and $\eR^1j_*\dZ/n$ is the
extension by zero of a locally free sheaf of rank one on $T$ (isomorphic,
via $d_2$, to $\dZ/n(-1)$). Since the map $a$ defined above is injective
(as may be checked fibre by fibre), it is an isomorphism, which proves
(\ref{I:5-3-4}).

\subsubsection{}\label{I:5-3-5}

We refer to \cite[XVI.4,5]{sga4} for proofs of the following applications
of the acyclicity theorem (\ref{I:5-2-1}).

\begin{theorem}\label{I:5-3-6} % originally 5.3.5.1
Let $f:X\to S$ be a morphism of schemes of finite type over $\dC$ and
let $\sF$ be a constructible sheaf on $X$. Then
\[
  \an{(\eR^q f_*\sF)} \sim \eR^q \an f_* (\an\sF)
\]
(cf. \ref{I:4-6-3}: for ordinary cohomology it is necessary to assume
that $\sF$ is constructible, not merely torsion).
\end{theorem}

\begin{theorem}\label{I:5-3-7} % originally 5.3.5.2
Let $f:X\to S$ be a morphism of schemes of finite type over a field $k$
of characteristic $0$, and let $\sF$ be a constructible sheaf on $X$.
Then the $\eR^qf_*\sF$ are also constructible.
\end{theorem}

The proof uses resolution of singularities and (\ref{I:5-3-4}). It is
generalized to a morphism of finite type between excellent schemes of
characteristic $0$ in \cite[XIX.5]{sga4}. Another proof, independent of
resolution, is given in this volume (\nameref{VII}, \ref{VII:1-1})
and applies to a morphism of schemes of finite type over a field or over
a Dedekind ring.

\section{Poincaré duality}

\subsection{Introduction}\label{I:6-1}

Let $X$ be an oriented topological manifold of pure dimension $N$, and
suppose that $X$ admits a finite open covering
$\cU=(U_i)_{1\leqslant i\leqslant K}$ such that every nonempty
intersection of the open subsets $U_i$ is homeomorphic to a ball. For
such a manifold, the Poincaré duality theorem can be presented as follows:
\begin{enumerate}[\indent A]
  \item The cohomology of $X$ is the \v Cech cohomology associated with
    the covering $\cU$. It is the cohomology of the complex
    \begin{equation}\label{I:eq:6-1}
      0 \to \dZ^{A_0} \to \dZ^{A_1} \to \cdots
    \end{equation}
    where
    $A_k=\{(i_0,\dotsc,i_k):i_0<\cdots<i_k\text{ and }U_{i_0}\cap\cdots\cap U_{i_k}\ne\varnothing\}$.
  \item For $a\in A_k$, $a=(i_0,\dotsc,i_k)$, let
    $U_a=U_{i_0}\cap\cdots\cap U_{i_k}$ and let $j_a$ be the inclusion
    of $U_a$ into $X$. The constant sheaf $\dZ$ on $X$ admits the
    (left) resolution
    \begin{equation}\label{I:eq:6-2}
      \xymatrix{
        \cdots \ar[r]
          & \displaystyle\bigoplus_{a\in A_1} j_{a!}\dZ \ar[r]
          & \displaystyle\bigoplus_{a\in A_0} j_{a!} \dZ \ar[r] \ar[d]
          & 0 \\
        & & \dZ \text{.}
      }
    \end{equation}
    The compactly supported cohomology $\h_c^\bullet(X,j_{a!}\dZ)$ is
    precisely the cohomology with proper support of the (oriented) ball
    $U_a$:
      \[
        \h_c^i(X,j_{a!}\dZ) = \begin{cases}
                                0   & \text{if $i\ne N$} \\
                                \dZ & \text{if $i = N$}
                              \end{cases}
      \]
\end{enumerate}

The hypercohomology spectral sequence for the complex
\eqref{I:eq:6-2} and cohomology with proper support therefore says that
$\h_c^i(X)$ is the $(i-N)$th cohomology group of the complex
\begin{equation}\label{I:eq:6-3}
  \cdots \to \dZ^{A_1} \to \dZ^{A_0} \to 0 \text{.}
\end{equation}
This complex is dual to the complex \eqref{I:eq:6-1}, whence Poincaré
duality.

The essential points of this construction are
\begin{enumerate}[\indent a)]
  \item the existence of a theory of cohomology with proper support;
  \item the fact that every point $x$ of a manifold $X$ of pure dimension
    $N$ has a fundamental system of open neighbourhoods $U$ for which
    \begin{equation}\label{I:eq:6-4}
      \h_c^i(U) = \begin{cases}
                    0   & \text{for $i\ne N$,} \\
                    \dZ & \text{for $i=N$.}
                  \end{cases}
    \end{equation}
\end{enumerate}

Poincaré duality in étale cohomology can be constructed on this model.
For $X$ smooth of pure dimension $N$ over an algebraically closed field
$k$, $n$ invertible on $X$, and $x$ a closed point of $X$, the key
point is to calculate the projective limit over the étale neighbourhoods
$U$ of $x$:
\begin{equation}\label{I:eq:6-5}
  \varprojlim \h_c^i(U,\dZ/n)
    = \begin{cases}
        0     & \text{if $i\ne 2N$} \\
        \dZ/n & \text{if $i=2N$.}
      \end{cases}
\end{equation}

Just as for topological manifolds one must first treat directly the case
of an open ball (or even simply the interval $(0,1)$), here one must
first treat directly the case of curves (\ref{I:6-2}). The local
acyclicity theorem for smooth morphisms then reduces
\eqref{I:eq:6-5} to this special case (\ref{I:6-3}).

The isomorphisms \eqref{I:eq:6-4} and \eqref{I:eq:6-5} are not
canonical: they depend on the choice of an \emph{orientation} of $X$.
For $n$ invertible on a scheme $X$, $\dmu_n$ is a sheaf of free
$\dZ/n$-modules of rank one. We denote its $N$th tensor power by
$\dZ/n(N)$ ($N\in\dZ$). The intrinsic form of the second line of
\eqref{I:eq:6-5} is
\begin{equation}\label{I:eq:6-5'}
  \varprojlim \h_c^{2N} \left(U,\dZ/n(N)\right) = \dZ/n \text{,}
\end{equation}
and $\dZ/n(N)$ is called the \emph{orientation sheaf} of $X$. Since the
sheaf $\dZ/n(N)$ is constant and isomorphic to $\dZ/n$, it may be taken
outside the limit sign (with $N$ fixed), and one may instead write
\begin{equation}\label{I:eq:6-5''}
  \varprojlim \h_c^{2N} (U,\dZ/n) = \dZ/n(-N) \text{,}
\end{equation}
and Poincaré duality will take the form of a perfect duality, with values
in $\dZ/n(-N)$, between $\h^i(X,\dZ/n)$ and
$\h_c^{2N-i}(X,\dZ/n)$.

\subsection{The case of curves}\label{I:6-2}

\subsubsection{}\label{I:6-2-1}

Let $\bar X$ be a smooth projective curve over an algebraically closed
field $k$, and let $n$ be invertible on $\bar X$. The proof of
(\ref{I:3-3-5}) gives, for connected $\bar X$, a canonical isomorphism
\[\xymatrix{
  \h^2(\bar X,\dmu_n) = \pic(\bar X)/n\pic(\bar X) \ar[r]^-{\deg}_-{\sim}
    & \dZ/n
}\]
Let $D$ be a reduced divisor of $\bar X$ and let
$X=\bar X\setminus D$:
\[\xymatrix{
  X \ar@{^{(}->}[r]^-j
    & \bar X
    & \ar@{_{(}->}[l]_-i D \text{.}
}\]
The exact sequence
$0\to j_!\dmu_n\to\dmu_n\to i_*\dmu_n\to0$ and the fact that
$\h^i(\bar X,i_*\dmu_n)=\h^i(D,\dmu_n)=0$ for $i>0$ give an
isomorphism
\[\xymatrix{
  \h_c^2(X,\dmu_n) = \h^2(\bar X,j_! \dmu_n) \ar[r]^-\sim
    & \h^2(\bar X,\dmu_n) \ar[r]^-\sim
    & \dZ/n \text{.}
}\]
For disconnected $\bar X$, similarly
\[
  \h_c^2(X,\dmu_n) = (\dZ/n)^{\pi_0(X)}
\]
and the \emph{trace morphism} is defined to be the sum
\[\xymatrix{
  \tr : \h_c^2(X,\dmu_n) \simeq (\dZ/n)^{\pi_0(X)} \ar[r]^-\Sigma
    & \dZ/n\text{.}
}\]

\begin{theorem}\label{I:6-2-2}
The pairing $\tr(a\smallsmile b)$ identifies each of the two groups
$\h^i(X,\dZ/n)$ and $\h_c^{2-i}(X,\dmu_n)$ with the dual (with values
in $\dZ/n$) of the other.
\end{theorem}

\emph{Transcendental proof.} Suppose that $\bar X$ is a smooth projective
curve over a trait $S$ and that $j:X\hookrightarrow\bar X$ is the
inclusion of the complement of a divisor $D$ étale over $S$. The
cohomology (resp. cohomology with proper support) of the special and
generic geometric fibres of $X/S$ are “the same,” i.e. the fibres of
locally constant sheaves on $S$. This follows from the analogous facts
for $\bar X$ and $D$, by means of the exact sequence
$0\to j_!\dZ/n\to\dZ/n\to{\dZ/n}_D\to0$ (for cohomology with proper
support) and the formulas $j_*\dZ/n=\dZ/n$,
$\eR^1j_*\dZ/n\simeq{\dZ/n}_D(-1)$, and
$\eR^ij_*\dZ/n=0$ ($i\geqslant2$) (for ordinary cohomology)
(\ref{I:5-3-4}).

This specialization principle reduces (\ref{I:6-2-2}) to the case in
which $k$ has characteristic $0$. By (\ref{I:5-3-3}), this case reduces
to $k=\dC$. Finally, when $k=\dC$, the groups
$\h^\bullet(X,\dZ/n)$ and $\h_c^\bullet(X,\dmu_n)$ agree with the
groups of the same name calculated for the classical topological space
$X_{\textnormal{cl}}$, and, through the isomorphism
$\dZ/n\to\dmu_n:x\mapsto\exp\left(\frac{2\pi ix}{n}\right)$, the trace
morphism is identified with “integration over the fundamental class.”
Thus (\ref{I:6-2-2}) follows from Poincaré duality for
$X_{\textnormal{cl}}$.

\subsubsection{Algebraic proof}\label{I:6-2-3}

For a very economical proof, see (\nameref{V}, \S\ref{V:2}).
Here is another, related to the autoduality of the Jacobian.

Return to the notation of (\ref{I:6-2-1}). We may suppose---and do
suppose---that $X$ is connected. The cases $i=0$ and $i=2$ being
trivial, we also suppose that $i=1$. Define $_D\dG_m$ by the exact
sequence $0\to_D\dG_m\to\dG_m\to i_*\dG_m\to0$ (sections of $\dG_m$
congruent to $1\mod D$). The group $\h^1(\bar X,_D\dG_m)$ classifies
invertible sheaves on $\bar X$ trivialized on $D$. It is the group of
points of $\pic_D(\bar X)$, an extension of $\dZ$ (the degree) by the
group of points of Rosenlicht's generalized Jacobian
$\pic_D^0(\bar X)$ (corresponding to conductor $1$ at each point of
$D$), itself an extension of the abelian variety $\pic^0(\bar X)$ by
the torus $\dG_m^D/(\text{diagonal $\dG_m$})$.

\begin{enumerate}[\indent a)]
  \item The exact sequence
    $0\to j_!\dmu_n\to_D\dG_m\xrightarrow{n}_D\dG_m\to0$ gives an
    isomorphism
    \begin{equation*}\tag{6.2.3.1}\label{I:eq:6-2-3-1}
      \h_c^1(X,\dmu_n) = \pic_D^0(\bar X)_n \text{.}
    \end{equation*}
  \item The map sending $x\in X(k)$ to the class of the invertible sheaf
    $\sO(x)$ on $\bar X$, trivialized by $1$ on $D$, comes from a
    morphism
    \[
      f:X \to \pic_D(\bar X) \text{.}
    \]
\end{enumerate}
For the remainder, fix a base point $0$ and put
$f_0(x)=f(x)-f(0)$. For every homomorphism
$v:\pic_D^0(\bar X)_n\to\dZ/n$, let
$\bar v\in\h^1\left(\pic_D^0(\bar X),\dZ/n\right)$ be the image under
$v$ of the class in
$\h^1\left(\pic_D^0(\bar X),\pic_D^0(\bar X)_n\right)$ of the torsor
defined by the extension
\[\xymatrix{
  0 \ar[r]
    & \pic_D^0(\bar X)_n \ar[r]
    & \pic_D^0(\bar X) \ar[r]^-n
    & \pic_D^0(\bar X) \ar[r]
    & 0
}\]
Geometric class field theory (as presented in Serre \cite{se59}) shows
that the map $v\mapsto f_0^*(\bar v)$
\begin{equation*}\tag{6.2.3.2}\label{I:eq:6-2-3-2}
  \hom\left(\pic_D^0(\bar X)_n,\dZ/n\right) \to \h^1(X,\dZ/n)
\end{equation*}
is an isomorphism. To deduce (\ref{I:6-2-2}) from
\eqref{I:eq:6-2-3-1} and \eqref{I:eq:6-2-3-2}, it remains to know that
\begin{equation*}\tag{2.3.3}\label{I:eq:6-2-3-3}
  \tr\left(u \smallsmile f_0^*(\bar v)\right) = -v(u) \text{.}
\end{equation*}


\subsection{The general case}\label{I:6-3}

Let $X$ be a smooth algebraic variety of pure dimension $N$ over an
algebraically closed field $k$. To state the Poincaré duality theorem,
one must first define the trace morphism
\[
  \tr : \h_c^{2N}\left(X,\dZ/n(N)\right) \to \dZ/n\text{.}
\]
The definition is a painful unwinding starting from the case of curves
\cite[XVIII \S 2]{sga4}. One then has:

\begin{theorem}\label{I:6-3-1}
The pairing $\tr(a\smallsmile b)$ identifies each of the groups
$\h_c^i\left(X,\dZ/n(N)\right)$ and $\h^{2N-i}(X,\dZ/n)$ with the dual
of the other.
\end{theorem}

Let $x\in X$ be a closed point and let $X_x$ be the strict localization
of $X$ at $x$. For $U$ ranging over the étale neighbourhoods of $x$,
put
\begin{equation*}\tag{6.3.1.1}\label{I:eq:6-3-1-1}
  \h_c^\bullet(X_x,\dZ/n) = \varprojlim \h_c^\bullet(U,\dZ/n) \text{.}
\end{equation*}
It would be preferable to consider instead the pro-object
$\text{“}\varprojlim\text{”}\h_c^\bullet(U,\dZ/n)$, but since the
groups involved are finite, the difference is inessential. As we tried
to explain in the introduction, (\ref{I:6-3-1}) follows from the facts
that
\begin{equation*}\tag{6.3.1.2}\label{I:eq:6-3-1-2}
\begin{aligned}
  \h_c^i(X_x,\dZ/n) &= 0 \text{ for $i\ne 2N$ and that } \\
  \tr : \h_c^{2 N}\left(X_x,  \dZ/n(N)\right) &\iso \dZ/n \text{ is an isomorphism.}
\end{aligned}
\end{equation*}

The case $N=0$ is trivial. If $N>0$, let $Y_y$ be the strict
localization at a closed point of a smooth scheme $Y$ of pure dimension
$N-1$, and let $f:X_x\to Y_y$ be an essentially smooth morphism (of
relative dimension one). The proof uses the Leray spectral sequence for
cohomology with proper support for $f$ to reduce to the case of curves.
Because the “cohomology with proper support” under consideration is
defined as limits \eqref{I:eq:6-3-1-1}, the existence of such a spectral
sequence raises various passage-to-the-limit problems, treated in too much
detail in \cite[XVIII]{sga4}. For every geometric point $z$ of $Y_y$,
one has
\[
  (\eR^if_! \dZ/n)_z = \h_c^i(f^{-1}(z),\dZ/n) \text{.}
\]
The geometric fibre $f^{-1}(z)$ is a projective limit of smooth curves
over an algebraically closed field. It satisfies Poincaré duality. Its
ordinary cohomology is given by the local acyclicity theorem for smooth
morphisms:
\[
  \h^i(f^{-1}(z),\dZ/n)
    = \begin{cases}
        \dZ/n & \text{for $i=0$} \\
        0     & \text{for $i>0$.}
      \end{cases}
\]
By duality,
\[
  \h_c^i(f^{-1}(z),\dZ/n)
    = \begin{cases}
        \dZ/n(-1) & \text{for $i=2$} \\
        0         & \text{for $i\ne 2$.}
      \end{cases}
\]
and the Leray spectral sequence becomes
\[
  \h_c^i\left(X_x,\dZ/n(N)\right) = \h_c^{i-2}\left(Y_y,\dZ/n(N-1)\right) \text{.}
\]
The result follows by induction on $N$.

\subsection{Variants and applications}\label{I:6-4}

One can construct in étale cohomology a “duality formalism”
($=$ functors $\eR f_*$, $\eR f_!$, $f^*$, and $\eR f^!$ satisfying various
compatibilities and adjunction formulas) parallel to the one that exists
in coherent cohomology. In this language, the results of the preceding
paragraph can be rewritten as follows: if $f:X\to S$ is smooth of pure
relative dimension $N$ and $S=\spec(k)$ with $k$ algebraically closed,
then
\[
  \eR f^! \dZ/n = \dZ/n [2N](N) \text{.}
\]
This result holds without any hypothesis on $S$. It has the following
corollary.

\begin{corollary}\label{I:6-4-1}
If $f$ is smooth of pure relative dimension $N$ and the sheaves
$\eR^if_!\dZ/n$ are locally constant, then the sheaves
$\eR^if_*\dZ/n$ are also locally constant, and
\[
  \eR^i f_* \dZ/n = \const\hom\left(\eR^{2 N-i} f_! \dZ/n(N),\dZ/n\right)\text{.}
\]
\end{corollary}

In particular, under the hypotheses of the corollary, the sheaves
$\eR^if_*\dZ/n$ are constructible. From this one can show that, if $S$
is of finite type over the spectrum of a field or of a Dedekind ring,
then, for every morphism of finite type $f:X\to S$ and every constructible
sheaf $\sF$ on $X$, the sheaves $\eR^if_*\sF$ are constructible
(\nameref{VII}, \ref{VII:1-1}).
